\Chapter{40}

\Verse{1}
{
    \zA{话说宝玉听了，忙进来看时，只见琥珀站在屏风跟前，说：“快去罢，立等你 说话呢。”}
}
{
    \eA{As soon as Pao-yü, we will now explain, heard what the lad told him, he
        rushed with eagerness inside. When he came to look about him, he discovered
        Hu Po standing in front of the screen. "Be quick and go," she urged.}
}
{
}

\Verse{2}
{
    \zA{宝玉来至上房，只见贾母正和王夫人众姐妹商议给史湘云还席。}
}
{
    \eA{"They're waiting to speak to you." Pao-yü wended his way into the drawing
        rooms. Here he found dowager lady Chia, consulting with Madame Wang and the
        whole body of young ladies,}
}
{
}

\Verse{3}
{
    \zA{宝玉因 说：“我有个主意：既没有外客，吃的东西也别定了样数，谁素日爱吃的，拣样儿 做几样。}
}
{
    \eA{about the return feast to be given to Shih Hsiang-yün. "I've got a plan to
        suggest," he consequently interposed. "As there are to be no outside guests,
        the eatables too should not be limited to any kind or number. A few of such
        dishes, as have ever been to the liking of any of us,}
}
{
}

\Verse{4}
{
    \zA{也不必按桌席，每人跟前摆一张高几，各人爱吃的东西一两样，再一个十 锦攒心盒子、自斟壶，岂不别致？”}
}
{
    \eA{should be fixed upon and prepared for the occasion. Neither should any
        banquet be spread, but a high teapoy can be placed in front of each, with
        one or two things to suit our particular tastes. Besides, a painted box with
        partitions and a decanter.}
}
{
}

\Verse{5}
{
    \zA{贾母听了，说：“很是。”}
}
{
    \eA{Won't this be an original way?" "Capital!" shouted old lady Chia.}
}
{
}

\Verse{6}
{
    \zA{即命人传与厨房： “明日就拣我们爱吃的东西做了，按着人数，再装了盒子来，早饭也摆在园里吃。”}
}
{
    \eA{"Go and tell the people in the cook house," she forthwith ordered a servant,
        "to get ready to-morrow such dishes as we relish, and to put them in as many
        boxes as there will be people, and bring them over. We can have breakfast
        too in the garden."}
}
{
}

\Verse{7}
{
    \zA{商议之间，早又掌灯，一夕无话。}
}
{
    \eA{But while they were deliberating, the time came to light the lamps.}
}
{
}

\Verse{8}
{
    \zA{次日清早起来，可喜这日天气清朗。}
}
{
    \eA{Nothing of any note transpired the whole night. The next day, they got up at
        early dawn.}
}
{
}

\Verse{9}
{
    \zA{李纨清晨起来，看着老婆子丫头们扫那些 落叶，并擦抹桌椅，预备茶酒器皿。}
}
{
    \eA{The weather, fortunately, was beautifully clear. Li Wan turned out of bed at
        daybreak. She was engaged in watching the old matrons and servant-girls
        sweeping the fallen leaves, rubbing the tables and chairs,}
}
{
}

\Verse{10}
{
    \zA{只见丰儿带了刘老老板儿进来，说：“大奶奶 倒忙的很。”}
}
{
    \eA{and preparing the tea and wine vessels, when she perceived Feng Erh usher in
        old goody Liu and Pan Erh. "You're very busy, our senior lady!"}
}
{
}

\Verse{11}
{
    \zA{李纨笑道：“我说你昨儿去不成，只忙着要去。”}
}
{
    \eA{they said. "I told you that you wouldn't manage to start yesterday," Li Wan
        smiled, "but you were in a hurry to get away."}
}
{
}

\Verse{12}
{
    \zA{刘老老笑道：“老 太太留下我，叫我也热闹一天去。”}
}
{
    \eA{"Your worthy old lady," goody Liu replied laughingly, "wouldn't let me go.
        She wanted me to enjoy myself too for a day before I went."}
}
{
}

\Verse{13}
{
    \zA{丰儿拿了几把大小钥匙，说道：“我们奶奶说 了，外头的高几儿怕不够使，不如开了楼，把那收的拿下来使一天罢。}
}
{
    \eA{Feng Erh then produced several large and small keys. "Our mistress Lien
        says," she remarked, "that she fears that the high teapoys which are out are
        not enough, and she thinks it would be as well to open the loft and take out
        those that are put away and use them for a day.}
}
{
}

\Verse{14}
{
    \zA{奶奶原该亲 自来，因和太太说话呢，请大奶奶开了，带着人搬罢。”}
}
{
    \eA{Our lady should really have come and seen to it in person, but as she has
        something to tell Madame Wang, she begs your ladyship to open the place,}
}
{
}

\Verse{15}
{
    \zA{李氏便命素云接了钥匙。}
}
{
    \eA{and get a few servants to bring them out." Li Wan there and then told Su Yün
        to take the keys.}
}
{
}

\Verse{16}
{
    \zA{又命婆子出去，把二门上小厮叫几个来。}
}
{
    \eA{She also bade a matron go out and call a few servant-boys from those on duty
        at the second gate.}
}
{
}

\Verse{17}
{
    \zA{李氏站在大观楼下往上看着，命人上去开 了缀锦阁，一张一张的往下抬。}
}
{
    \eA{When they came, Li Wan remained in the lower story of the Ta Kuan loft, and
        looking up, she ordered the servants to go and open the Cho Chin hall and to
        bring the teapoys one by one.}
}
{
}

\Verse{18}
{
    \zA{小厮、老婆子、丫头一齐动手，抬了二十多张下来。}
}
{
    \eA{The young servant-lads, matrons and servant-maids then set to work, in a
        body, and carried down over twenty of them.}
}
{
}

\Verse{19}
{
    \zA{李纨道：“好生着，别慌慌张张鬼赶着似的，仔细碰了牙子！”}
}
{
    \eA{"Be careful with them," shouted Li Wan. "Don't be bustling about just as if
        you were being pursued by ghosts! Mind you don't break the tenons!"}
}
{
}

\Verse{20}
{
    \zA{又回头向刘老老笑 道：“老老也上去瞧瞧。”}
}
{
    \eA{Turning her head round, "old dame," she observed, addressing herself
        smilingly to goody Liu, "go upstairs too and have a look!"}
}
{
}

\Verse{21}
{
    \zA{刘老老听说巴不得一声儿，拉了板儿登梯上去。}
}
{
    \eA{Old goody Liu was longing to satisfy her curiosity, so at the bare mention
        of the permission, she uttered just one word ("come") and,}
}
{
}

\Verse{22}
{
    \zA{进里面 只见乌压压的堆着些围屏桌椅、大小花灯之类，虽不大认得，只见五彩？}
}
{
    \eA{dragging Pan Erh along, she trudged up the stairs. On her arrival inside,
        she espied, pile upon pile, a whole heap of screens, tables and chairs,
        painted lanterns of different sizes, and other similar articles.}
}
{
}

\Verse{23}
{
    \zA{灼，各有 奇妙，念了几声佛便下来了。}
}
{
    \eA{She could not, it is true, make out the use of the various things, but, at
        the sight of so many colours,}
}
{
}

\Verse{24}
{
    \zA{然后锁上门，一齐下来。}
}
{
    \eA{of such finery and of the unusual beauty of each article, she muttered time
        after time the name of Buddha,}
}
{
}

\Verse{25}
{
    \zA{李纨道：“恐怕老太太高兴， 越发把船上划子、篙、桨、遮阳幔子，都搬下来预备着。”}
}
{
    \eA{and then forthwith wended her way downstairs. Subsequently (the servants)
        locked the doors and every one of them came down. "I fancy," cried Li Wan,
        "that our dowager lady will feel disposed (to go on the water),}
}
{
}

\Verse{26}
{
    \zA{众人答应，又复开了门， 色色的搬下来。}
}
{
    \eA{so you'd better also get the poles, oars and awnings for the boats and keep
        them in readiness."}
}
{
}

\Verse{27}
{
    \zA{命小厮传驾娘们，到船坞里撑出两只船来。}
}
{
    \eA{The servants expressed their obedience. Once more they unlocked the doors,
        and carried down everything required.}
}
{
}

\Verse{28}
{
    \zA{正乱着，只见贾母已带了一群人进来了，李纨忙迎上去，笑道：“老太太高兴， 倒进来了；我只当还没梳头呢，才掐了菊花要送去。”}
}
{
    \eA{She then bade a lad notify the boatwomen go to the dock and punt out two
        boats. But while all this bustle was going on, they discovered that dowager
        lady Chia had already arrived at the head of a whole company of people. Li
        Wan promptly went up to greet them. "Dear venerable senior," she smiled,
        "you must be in good spirits to have come in here!}
}
{
}

\Verse{29}
{
    \zA{一面说，一面碧月早已捧过 一个大荷叶式的翡翠盘子来，里面养着各色折枝菊花。}
}
{
    \eA{Imagining that you hadn't as yet combed your hair, I just plucked a few
        chrysanthemums, meaning to send them to you." While she spoke, Pi Yüeh at
        once presented to her a jadite tray,}
}
{
}

\Verse{30}
{
    \zA{贾母便拣了一朵大红的簪在 鬓上，因回头看见了刘老老，忙笑道：“过来带花儿。”}
}
{
    \eA{of the size of a lotus leaf, containing twigs cut from every species of
        chrysanthemum. Old lady Chia selected a cluster of deep red and pinned it in
        her hair about her temples. But turning round, she noticed old goody Liu.}
}
{
}

\Verse{31}
{
    \zA{一语未完，凤姐儿便拉过 刘老老来，笑道：“让我打扮你。”}
}
{
    \eA{"Come over here," she vehemently cried with a smile; "and put on a few
        flowers." Scarcely was this remark concluded, than lady Feng dragged goody
        Liu forward.}
}
{
}

\Verse{32}
{
    \zA{说着，把一盘子花，横三竖四的插了一头。}
}
{
    \eA{"Let me deck you up!" she laughed. With these words, she seized a whole
        plateful of flowers and stuck them three this way,}
}
{
}

\Verse{33}
{
    \zA{贾 母和众人笑的了不得。}
}
{
    \eA{four that way, all over her head. Old lady Chia, and the whole party were
        greatly amused;}
}
{
}

\Verse{34}
{
    \zA{刘老老也笑道：“我这头也不知修了什么福，今儿这样体面 起来。”}
}
{
    \eA{so much so, that they could not check themselves. "I wonder," shouted goody
        Liu smiling, "what blessings I have brought upon my head that such honours
        are conferred upon it to-day!"}
}
{
}

\Verse{35}
{
    \zA{众人笑道：“你还不拔下来摔到他脸上呢，把你打扮的成了老妖精了。”}
}
{
    \eA{"Don't you yet pull them away," they all laughed, "and chuck them in her
        face! She has got you up in such a way as to make a regular old elf of you!"
        "I'm an old hag, I admit,"}
}
{
}

\Verse{36}
{
    \zA{刘老老笑道：“我虽老了，年轻时也风流，爱个花儿粉儿的，今儿索性作个老风流！”}
}
{
    \eA{goody Liu pursued with a laugh; "but when I was young, I too was pretty and
        fond of flowers and powder! But the best thing I can do now is to keep to
        such fineries as befit my advanced age!" While they bandied words,}
}
{
}

\Verse{37}
{
    \zA{说话间，已来至沁芳亭上，丫鬟们抱了个大锦褥子来，铺在栏杆榻板上。}
}
{
    \eA{they reached the Hsin Fang pavilion. The waiting maids brought a large
        embroidered rug and spread it over the planks of the divan near the
        balustrade. On this rug dowager lady Chia sat,}
}
{
}

\Verse{38}
{
    \zA{贾母 倚栏坐下，命刘老老也坐在旁边，因问他：“这园子好不好？”}
}
{
    \eA{with her back leaning against the railing; and, inviting goody Liu to also
        take a seat next to her, "Is this garden nice or not?" she asked her. Old
        goody Liu invoked Buddha several times.}
}
{
}

\Verse{39}
{
    \zA{刘老老念佛说道： “我们乡下人，到了年下，都上城来买画儿贴。}
}
{
    \eA{"We country-people," she rejoined, "do invariably come, at the close of each
        year, into the city and buy pictures and stick them about.}
}
{
}

\Verse{40}
{
    \zA{闲了的时候儿大家都说：‘怎么得 到画儿上逛逛！’}
}
{
    \eA{And frequently do we find ourselves in our leisure moments wondering how we
        too could manage to get into the pictures, and walk about the scenes they
        represent.}
}
{
}

\Verse{41}
{
    \zA{想着画儿也不过是假的，那里有这个真地方儿？}
}
{
    \eA{I presumed that those pictures were purely and simply fictitious, for how
        could there be any such places in reality?}
}
{
}

\Verse{42}
{
    \zA{谁知今儿进这园 里一瞧，竟比画儿还强十倍!怎么得有人也照着这个园子画一张，我带了家去给他 们见见，死了也得好处。”}
}
{
    \eA{But, contrary to my expectations, I found, as soon as I entered this garden
        to-day and had a look about it, that it was, after all, a hundred times
        better than these very pictures. But if only I could get some one to make me
        a sketch of this garden, to take home with me and let them see it,}
}
{
}

\Verse{43}
{
    \zA{贾母听说，指着惜春笑道：“你瞧我这个小孙女儿，他 就会画，等明儿叫他画一张如何？”}
}
{
    \eA{so that when we die we may have reaped some benefit!" Upon catching the wish
        she expressed, dowager lady Chia pointed at Hsi Ch'un. "Look at that young
        granddaughter of mine!" she smiled. "She's got the knack of drawing.}
}
{
}

\Verse{44}
{
    \zA{刘老老听了，喜的忙跑过来拉着惜春，说道： “我的姑娘!你这么大年纪儿，又这么个好模样儿，还有这个能干，别是个神仙托 生的罢？”}
}
{
    \eA{So what do you say to my asking her to-morrow to make a picture for you?"
        This suggestion filled goody Liu with enthusiasm and speedily crossing over,
        she clasped Hsi Ch'un in her arms. "My dear Miss!" she cried, "so young in
        years, and yet so pretty,}
}
{
}

\Verse{45}
{
    \zA{贾母众人都笑了。}
}
{
    \eA{and so accomplished too! Mightn't you be a spirit come to life!"}
}
{
}

\Verse{46}
{
    \zA{歇了歇，又领着刘老老都见识见识。}
}
{
    \eA{After old lady Chia had had a little rest, she in person took goody Liu and
        showed her everything there was to be seen.}
}
{
}

\Verse{47}
{
    \zA{先到了潇湘馆。}
}
{
    \eA{First, they visited the Hsiao Hsiang lodge.}
}
{
}

\Verse{48}
{
    \zA{一进门，只见两边翠竹夹 路，土地下苍苔布满，中间羊肠一条石子漫的甬路。}
}
{
    \eA{The moment they stepped into the entrance, a narrow avenue, flanked on
        either side with kingfisher-like green bamboos, met their gaze. The earth
        below was turfed all over with moss. In the centre, extended a tortuous
        road,}
}
{
}

\Verse{49}
{
    \zA{刘老老让出来与贾母众人走， 自己却走土地。}
}
{
    \eA{paved with pebbles. Goody Liu left dowager lady Chia and the party walk on
        the raised road, while she herself stepped on the earth.}
}
{
}

\Verse{50}
{
    \zA{琥珀拉他道：“老老你上来走，看青苔滑倒了。”}
}
{
    \eA{But Hu Po tugged at her. "Come up, old dame, and walk here!" she exclaimed.
        "Mind the fresh moss is slippery and you might fall."}
}
{
}

\Verse{51}
{
    \zA{刘老老道：“不 相干，我们走熟了，姑娘们只管走罢。}
}
{
    \eA{"I don't mind it!" answered goody Liu. "We people are accustomed to walking
        (on such slippery things)! So, young ladies, please proceed.}
}
{
}

\Verse{52}
{
    \zA{可惜你们的那鞋，别沾了泥。”}
}
{
    \eA{And do look after your embroidered shoes! Don't splash them with mud."}
}
{
}

\Verse{53}
{
    \zA{他只顾上头 和人说话，不防脚底下果踩滑了，“咕咚”一交跌倒，众人都拍手呵呵的大笑。}
}
{
    \eA{But while bent upon talking with those who kept on the raised road, she
        unawares reached a spot, which was actually slippery, and with a sound of
        "ku tang" she tumbled over.}
}
{
}

\Verse{54}
{
    \zA{贾 母笑骂道：“小蹄子们，还不搀起来，只站着笑！”}
}
{
    \eA{The whole company clapped their hands and laughed boisterously. "You young
        wenches," shouted out dowager lady Chia,}
}
{
}

\Verse{55}
{
    \zA{说话时，刘老老已爬起来了， 自己也笑了，说道：“才说嘴，就打了嘴了。”}
}
{
    \eA{"don't you yet raise her up, but stand by giggling?" This reprimand was
        still being uttered when goody Liu had already crawled up. She too was
        highly amused.}
}
{
}

\Verse{56}
{
    \zA{贾母问他：“可扭了腰了没有？}
}
{
    \eA{"Just as my mouth was bragging," she observed, "I got a whack on the lips!"}
}
{
}

\Verse{57}
{
    \zA{叫 丫头们捶捶。”}
}
{
    \eA{"Have you perchance twisted your waist?"}
}
{
}

\Verse{58}
{
    \zA{刘老老道：“那里说的我这么娇嫩了？}
}
{
    \eA{inquired old lady Chia. "Tell the servant-girls to pat it for you!" "What an
        idea!" retorted goody Liu,}
}
{
}

\Verse{59}
{
    \zA{那一天不跌两下子？}
}
{
    \eA{"am I so delicate? What day ever goes by without my tumbling down a couple
        of times?}
}
{
}

\Verse{60}
{
    \zA{都要捶起 来，还了得呢。”}
}
{
    \eA{And if I had to be patted every time wouldn't it be dreadful!"}
}
{
}

\Verse{61}
{
    \zA{紫鹃早打起湘帘，贾母等进来坐下。}
}
{
    \eA{Tzu Chuan had at an early period raised the speckled bamboo portiere.
        Dowager lady Chia and her companions entered and seated themselves.}
}
{
}

\Verse{62}
{
    \zA{黛玉亲自用小茶盘儿捧了一盖碗茶来奉与 贾母。}
}
{
    \eA{Lin Tai-yü with her own hands took a small tray and came to present a
        covered cup of tea to her grandmother. "We won't have any tea!"}
}
{
}

\Verse{63}
{
    \zA{王夫人道：“我们不吃茶，姑娘不用倒了。”}
}
{
    \eA{Madame Wang interposed, "so, miss, you needn't pour any." Lin Tai-yü,
        hearing this,}
}
{
}

\Verse{64}
{
    \zA{黛玉听说，便命丫头把自己窗 下常坐的一张椅子挪到下手，请王夫人坐了。}
}
{
    \eA{bade a waiting-maid fetch the chair from under the window where she herself
        often sat, and moving it to the lower side, she pressed Madame Wang into it.
        But goody Liu caught sight of the pencils and inkslabs,}
}
{
}

\Verse{65}
{
    \zA{刘老老因见窗下案上设着笔砚，又见 书架上放着满满的书，刘老老道：“这必定是那一位哥儿的书房了？”}
}
{
    \eA{lying on the table placed next to the window, and espied the bookcase piled
        up to the utmost with books. "This must surely," the old dame ejaculated,
        "be some young gentleman's study!" "This is the room of this
        granddaughter-in-law of mine,"}
}
{
}

\Verse{66}
{
    \zA{贾母笑指黛 玉道：“这是我这外孙女儿的屋子。”}
}
{
    \eA{dowager lady Chia explained, smilingly pointing to Tai-yü. Goody Liu
        scrutinised Lin Tai-yü with intentness for a while.}
}
{
}

\Verse{67}
{
    \zA{刘老老留神打量了黛玉一番，方笑道：“这 那里像个小姐的绣房？}
}
{
    \eA{"Is this anything like a young lady's private room?" she then observed with
        a smile. "Why, in very deed, it's superior to any first class library!"}
}
{
}

\Verse{68}
{
    \zA{竟比那上等的书房还好呢。”}
}
{
    \eA{"How is it I don't see Pao-yü?" his grandmother Chia went on to inquire.}
}
{
}

\Verse{69}
{
    \zA{贾母因问：“宝玉怎么不见？”}
}
{
    \eA{"He's in the boat, on the pond," the waiting-maids, with one voice,}
}
{
}

\Verse{70}
{
    \zA{众丫头们答说：“在池子里船上呢。”}
}
{
    \eA{returned for answer. "Who also got the boats ready?" old lady Chia asked.
        "The loft was open just now so they were taken out,"}
}
{
}

\Verse{71}
{
    \zA{贾母道：“谁又预备下船了？”}
}
{
    \eA{Li Wan said, "and as I thought that you might, venerable senior, feel
        inclined to have a row,}
}
{
}

\Verse{72}
{
    \zA{李纨忙回说： “才开楼拿的。}
}
{
    \eA{I got everything ready." After listening to this explanation,}
}
{
}

\Verse{73}
{
    \zA{我恐怕老太太高兴，就预备下了。”}
}
{
    \eA{dowager lady Chia was about to pass some remark, but some one came and
        reported to her that Mrs. Hsüeh had arrived.}
}
{
}

\Verse{74}
{
    \zA{贾母听了，方欲说话时，有人回说：“姨太太来了。”}
}
{
    \eA{No sooner had old lady Chia and the others sprung to their feet than they
        noticed that Mrs. Hsüeh had already made her appearance.}
}
{
}

\Verse{75}
{
    \zA{贾母等刚站起来，只见 薛姨妈早进来了，一面归坐，笑道：“今儿老太太高兴，这早晚就来了。”}
}
{
    \eA{While taking a seat: "Your venerable ladyship," she smiled, "must be in
        capital spirits to-day to have come at this early hour!" "It's only this
        very minute that I proposed that any one who came late, should be fined,"
        dowager lady Chia laughed,}
}
{
}

\Verse{76}
{
    \zA{贾母笑 道：“我才说，来迟了的要罚他，不想姨太太就来迟了。”}
}
{
    \eA{"and, who'd have thought it, here you, Mrs. Hsüeh, arrive late!" After they
        had indulged in good-humoured raillery for a time, old lady Chia's attention
        was attracted by the faded colour of the gauze on the windows,}
}
{
}

\Verse{77}
{
    \zA{说笑一回。}
}
{
    \eA{and she addressed herself to Madame Wang.}
}
{
}

\Verse{78}
{
    \zA{贾母因见窗 上纱颜色旧了，便和王夫人说道：“这个纱新糊上好看，过了后儿就不翠了。}
}
{
    \eA{"This gauze," she said, "may have been nice enough when it was newly pasted,
        but after a time nothing remained of kingfisher green. In this court too
        there are no peach or apricot trees and these bamboos already are green in
        themselves, so were this shade of green gauze to be put up again,}
}
{
}

\Verse{79}
{
    \zA{这院 子里头又没有个桃杏树，这竹子已是绿的，再拿绿纱糊上，反倒不配。}
}
{
    \eA{it would, instead of improving matters, not harmonise with the surroundings.
        I remember that we had at one time four or five kinds of coloured gauzes for
        sticking on windows, so give her some to-morrow to change that on there."}
}
{
}

\Verse{80}
{
    \zA{我记得咱们 先有四五样颜色糊窗的纱呢。}
}
{
    \eA{"When I opened the store yesterday," hastily put in Lady Feng, "I noticed
        that there were still in those boxes,}
}
{
}

\Verse{81}
{
    \zA{明儿给他把这窗上的换了。”}
}
{
    \eA{made of large planks, several rolls of 'cicada wing' gauze of silvery red
        colour.}
}
{
}

\Verse{82}
{
    \zA{凤姐儿忙道：“昨儿我 开库房，看见大板箱里还有好几匹银红蝉翼纱，也有各样折枝花样的，也有‘流云
        蝙蝠’花样的，也有‘百蝶穿花’花样的，颜色又鲜，纱又轻软，我竟没见这个样 的，拿了两匹出来，做两床绵纱被，想来一定是好的。”}
}
{
    \eA{There were also several rolls with designs of twigs of flowers of every
        kind, several with 'the rolling clouds and bats' pattern, and several with
        figures representing hundreds of butterflies, interspersed among flowers.
        The colours of all these were fresh, and the gauze supple. But I failed to
        see anything of the kind you speak of. Were two rolls taken (from those I
        referred to), and a couple of bed-covers of embroidered gauze made out of
        them,}
}
{
}
\ChapterImage{img/chapters/095第40回 史太君兩宴大觀園 劉姥姥初遊大觀園.jpg}


\Verse{83}
{
    \zA{贾母听了笑道：“呸，人 人都说你没有没经过没见过的，连这个纱还不能认得，明儿还说嘴。”}
}
{
    \eA{they would, I fancy, be a pretty sight!" "Pshaw!" laughed old lady Chia,
        "every one says that there's nothing you haven't gone through and nothing
        you haven't seen, and don't you even know what this gauze is? Will you again
        brag by and bye,}
}
{
}

\Verse{84}
{
    \zA{薛姨妈等都 笑说：“凭他怎么经过见过，怎么敢比老太太呢!老太太何不教导了他，连我们也 听听。”}
}
{
    \eA{after this?" Mrs. Hsüeh and all the others smiled. "She may have gone
        through a good deal," they remarked, "but how can she ever presume to pit
        herself against an old lady like you? So why don't you, venerable senior,
        tell her what it is so that we too may be edified."}
}
{
}

\Verse{85}
{
    \zA{凤姐儿也笑说：“好祖宗，教给我罢。”}
}
{
    \eA{Lady Feng too gave a smile. "My dear ancestor," she pleaded, "do tell me
        what it is like."}
}
{
}

\Verse{86}
{
    \zA{贾母笑向薛姨妈众人道：“那个 纱，比你们的年纪还大呢，怪不得他认做蝉翼纱，原也有些像。}
}
{
    \eA{Dowager lady Chia thereupon proceeded to enlighten Mrs. Hsüeh and the whole
        company. "That gauze is older in years than any one of you," she said. "It
        isn't therefore to be wondered, if you make a mistake and take it for
        'cicada wing' gauze.}
}
{
}

\Verse{87}
{
    \zA{不知道的都认做蝉 翼纱。}
}
{
    \eA{But it really bears some resemblance to it; so much so, indeed, that any
        one,}
}
{
}

\Verse{88}
{
    \zA{正经名字叫‘软烟罗’。”}
}
{
    \eA{not knowing the difference, would imagine it to be the 'cicada wing' gauze.}
}
{
}

\Verse{89}
{
    \zA{凤姐儿道：“这个名儿也好听，只是我这么大了， 纱罗也见过几百样，从没听见过这个名色。”}
}
{
    \eA{Its true name, however, is 'soft smoke' silk." "This is also a nice sounding
        name," lady Feng agreed. "But up to the age I've reached, I have never heard
        of any such designation, in spite of the many hundreds of specimens of
        gauzes and silks,}
}
{
}

\Verse{90}
{
    \zA{贾母笑道：“你能活了多大？}
}
{
    \eA{I've seen." "How long can you have lived?" old lady Chia added smilingly,}
}
{
}

\Verse{91}
{
    \zA{见过几 样东西？}
}
{
    \eA{"and how many kinds of things can you have met,}
}
{
}

\Verse{92}
{
    \zA{就说嘴来了。}
}
{
    \eA{that you indulge in this tall talk?}
}
{
}

\Verse{93}
{
    \zA{那个软烟罗只有四样颜色：一样雨过天青，一样秋香色，一 样松绿的，一样就是银红的。}
}
{
    \eA{Of this 'soft smoke' silk, there only exist four kinds of colours. The one
        is red-blue; the other is russet; the other pine-green; the other
        silvery-red; and it's because, when made into curtains or stuck on
        window-frames,}
}
{
}

\Verse{94}
{
    \zA{要是做了帐子，糊了窗屉，远远的看着就和烟雾一样， 所以叫做‘软烟罗’。}
}
{
    \eA{it looks from far like smoke or mist, that it is called 'soft smoke' silk.
        The silvery-red is also called 'russet shadow' gauze. Among the gauzes used
        in the present day, in the palace above, there are none so supple and rich,}
}
{
}

\Verse{95}
{
    \zA{那银红的又叫做‘霞影纱’。}
}
{
    \eA{light and closely-woven as this!" "Not to speak of that girl Feng not having
        seen it,"}
}
{
}

\Verse{96}
{
    \zA{如今上用的府纱也没有这样软 厚轻密的了。”}
}
{
    \eA{Mrs. Hsüeh laughed, "why, even I have never so much as heard anything of
        it." While the conversation proceeded in this strain,}
}
{
}

\Verse{97}
{
    \zA{薛姨妈笑道：“别说凤丫头没见，连我也没听见过。”}
}
{
    \eA{lady Feng soon directed a servant to fetch a roll. "Now isn't this the
        kind!" dowager lady Chia exclaimed. "At first, we simply had it stuck on the
        window frames,}
}
{
}

\Verse{98}
{
    \zA{凤姐儿一面 说话，早命人取了一匹来了，贾母说：“可不是这个!先时原不过是糊窗屉，后来 我们拿这个做被做帐子试试，也竟好。}
}
{
    \eA{but we subsequently used it for covers and curtains, just for a trial, and
        really they were splendid! So you had better to-morrow try and find several
        rolls, and take some of the silvery-red one and have it fixed on the windows
        for her." While lady Feng promised to attend to her commission, the party
        scrutinised it,}
}
{
}

\Verse{99}
{
    \zA{明日就找出几匹来，拿银红的替他糊窗户。”}
}
{
    \eA{and unanimously extolled it with effusion. Old goody Liu too strained her
        eyes and examined it, and her lips incessantly muttered Buddha's name.}
}
{
}

\Verse{100}
{
    \zA{凤姐答应着。}
}
{
    \eA{"We couldn't," she ventured,}
}
{
}

\Verse{101}
{
    \zA{众人看了，都称赞不已。}
}
{
    \eA{"afford to make clothes of such stuff, much though we may long to do so;}
}
{
}

\Verse{102}
{
    \zA{刘老老也觑着眼看，口里不住的念佛，说道： “我们想做衣裳也不能，拿着糊窗子岂不可惜？”}
}
{
    \eA{and won't it be a pity to use it for sticking on windows?" "But it doesn't,
        after all, look well, when made into clothes," old lady Chia explained. Lady
        Feng hastily pulled out the lapel of the deep-red brocaded gauze jacket she
        had on,}
}
{
}

\Verse{103}
{
    \zA{贾母道：“倒是做衣裳不好看。”}
}
{
    \eA{and, facing dowager lady Chia and Mrs. Hsüeh, "Look at this jacket of mine,"}
}
{
}

\Verse{104}
{
    \zA{凤姐忙把自己身上穿的一件大红棉纱袄的襟子拉出来，向贾母薛姨妈道：“看我的 这袄儿。”}
}
{
    \eA{she remarked. "This is also of first-rate quality!" old lady Chia and Mrs.
        Hsüeh rejoined. "This is nowadays made in the palace for imperial use, but
        it can't possibly come up to this!"}
}
{
}

\Verse{105}
{
    \zA{贾母薛姨妈都说：“这也是上好的了，这是如今上用内造的，竟比不上 这个。”}
}
{
    \eA{"It's such thin stuff," lady Feng observed, "and do you still say that it
        was made in the palace for imperial use? Why, it doesn't, in fact, compare
        favourably with even this,}
}
{
}

\Verse{106}
{
    \zA{凤姐儿道：“这个薄片子还说是内造上用呢，竟连这个官用的也比不上啊。”}
}
{
    \eA{which is worn by officials!" "You'd better search again!" old lady Chia
        urged; "I believe there must be more of it! If there be, bring it all out,
        and give this old relative Liu a couple of rolls!}
}
{
}

\Verse{107}
{
    \zA{贾母道：“再找一找，只怕还有，要有就都拿出来，送这刘亲家两匹。}
}
{
    \eA{Should there be any red-blue, I'll make a curtain to hang up. What remains
        can be matched with some lining, and cut into a few double waistcoats for
        the waiting-maids to wear.}
}
{
}

\Verse{108}
{
    \zA{有雨过天青 的，我做一个帐子挂上。}
}
{
    \eA{It would be sheer waste to keep these things, as they will be spoilt by the
        damp." Lady Feng vehemently acquiesced;}
}
{
}

\Verse{109}
{
    \zA{剩的配上里子，做些个夹坎肩儿给丫头们穿，白收着霉坏 了。”}
}
{
    \eA{after which, she told a servant to take the gauze away. "These rooms are so
        small!" dowager lady Chia then observed, smiling. "We had better go
        elsewhere for a stroll."}
}
{
}

\Verse{110}
{
    \zA{凤姐儿忙答应了，仍命人送去。}
}
{
    \eA{"Every one says," old goody Liu put in, "that big people live in big houses!}
}
{
}

\Verse{111}
{
    \zA{贾母便笑道：“这屋里窄，再往别处逛去罢。”}
}
{
    \eA{When I saw yesterday your main apartments, dowager lady, with all those
        large boxes, immense presses, big tables,}
}
{
}

\Verse{112}
{
    \zA{刘老老笑道：“人人都说：‘大 家子住大房。’}
}
{
    \eA{and spacious beds to match, they did, indeed, present an imposing sight!
        Those presses are larger than our whole house;}
}
{
}

\Verse{113}
{
    \zA{昨儿见了老太太正房，配上大箱、大柜、大桌子、大床，果然威武。}
}
{
    \eA{yea loftier too! But strange to say there were ladders in the back court.
        'They don't also,' I thought, 'go up to the house tops to sun things, so
        what can they keep those ladders in readiness for?'}
}
{
}

\Verse{114}
{
    \zA{那柜子比我们一间房子还大还高。}
}
{
    \eA{Well, after that, I remembered that they must be required for opening the
        presses to take out or put in things. And that without those ladders,}
}
{
}

\Verse{115}
{
    \zA{怪道后院子里有个梯子，我想又不上房晒东西， 预备这梯子做什么？}
}
{
    \eA{how could one ever reach that height? But now that I've also seen these
        small rooms, more luxuriously got up than the large ones, and full of
        various articles,}
}
{
}

\Verse{116}
{
    \zA{后来我想起来，一定是为开顶柜取东西，离了那梯子怎么上得 去呢？}
}
{
    \eA{all so fascinating and hardly even known to me by name, I feel, the more I
        feast my eyes on them, the more unable to tear myself away from them."
        "There are other things still better than this,"}
}
{
}

\Verse{117}
{
    \zA{如今又见了这小屋子，更比大的越发齐整了。}
}
{
    \eA{lady Feng added. "I'll take you to see them all!" Saying this, they
        straightway left the Hsiao Hsiang lodge.}
}
{
}

\Verse{118}
{
    \zA{满屋里东西都只好看，可不知 叫什么。}
}
{
    \eA{From a distance, they spied a whole crowd of people punting the boats in the
        lake. "As they've got the boats ready,"}
}
{
}

\Verse{119}
{
    \zA{我越看越舍不得离了这里了！”}
}
{
    \eA{old lady Chia proposed, "we may as well go and have a row in them!" As she
        uttered this suggestion,}
}
{
}

\Verse{120}
{
    \zA{凤姐道：“还有好的呢，我都带你去瞧瞧。”}
}
{
    \eA{they wended their steps along the persicary-covered bank of the Purple Lily
        Isle. But before reaching the lake,}
}
{
}

\Verse{121}
{
    \zA{说着，一径离了潇湘馆，远远望见池中一群人在那里撑船。}
}
{
    \eA{they perceived several matrons advancing that way with large multi-coloured
        boxes in their hands, made all alike of twisted wire and inlaid with gold.}
}
{
}

\Verse{122}
{
    \zA{贾母道：“他们既 备下船，咱们就坐一回。”}
}
{
    \eA{Lady Feng hastened to inquire of Madame Wang where breakfast was to be
        served. "Ask our venerable senior,"}
}
{
}

\Verse{123}
{
    \zA{说着，向紫菱洲蓼溆一带走来。}
}
{
    \eA{Madame Wang replied, "and let them lay it wherever she pleases." Old lady
        Chia overheard her answer,}
}
{
}

\Verse{124}
{
    \zA{未至池前，只见几个婆 子手里都捧着一色摄丝戗金五彩大盒子走来，凤姐忙问王夫人：“早饭在那里摆？”}
}
{
    \eA{and turning her head round: "Miss Tertia," she said, "take the servants, and
        make them lay breakfast wherever you think best! We'll get into the boats
        from here." Upon catching her senior's wishes,}
}
{
}

\Verse{125}
{
    \zA{王夫人道：“问老太太在那里就在那里罢了。”}
}
{
    \eA{lady Feng retraced her footsteps, and accompanied by Li Wan, T'an Ch'un,
        Yüan Yang and Hu Po, she led off the servants,}
}
{
}

\Verse{126}
{
    \zA{贾母听说，便回头说：“你三妹妹 那里好，你就带了人摆去，我们从这里坐了船去。”}
}
{
    \eA{carrying the eatables, and other domestics, and came by the nearest way, to
        the Ch'iu Shuang library, where they arranged the tables in the Hsiao Ts'ui
        hall.}
}
{
}

\Verse{127}
{
    \zA{凤姐儿听说，便回身和李纨、 探春、鸳鸯、琥珀带着端饭的人等，抄着近路到了秋爽斋，就在晓翠堂上调开桌案。}
}
{
    \eA{"We daily say that whenever the gentlemen outside have anything to drink or
        eat, they invariably have some one who can raise a laugh and whom they can
        chaff for fun's sake," Yuan Yang smiled, "so let's also to-day get a female
        family-companion." Li Wan, being a person full of kindly feelings,}
}
{
}

\Verse{128}
{
    \zA{鸳鸯笑道：“天天咱们说外头老爷们吃酒吃饭，都有个凑趣儿的，拿他取笑儿。}
}
{
    \eA{did not fathom the insinuation, though it did not escape her ear. Lady Feng,
        however, thoroughly understood that she alluded to old goody Liu. "Let us
        too to-day," she smilingly remarked,}
}
{
}

\Verse{129}
{
    \zA{咱 们今儿也得了个女清客了。”}
}
{
    \eA{"chaff her for a bit of fun!" These two then began to mature their plans.}
}
{
}

\Verse{130}
{
    \zA{李纨是个厚道人，倒不理会；凤姐儿却听着是说刘老 老，便笑道：“咱们今儿就拿他取个笑儿。”}
}
{
    \eA{Li Wan chided them with a smile. "You people," she said, "don't know even
        how to perform the least good act! But you're not small children any more,
        and are you still up to these pranks? Mind, our venerable ancestor might
        call you to task!"}
}
{
}

\Verse{131}
{
    \zA{二人便如此这般商议。}
}
{
    \eA{"That has nothing whatever to do with you, senior lady," Yüan Yang laughed,}
}
{
}

\Verse{132}
{
    \zA{李纨笑劝道： “你们一点好事儿不做。}
}
{
    \eA{"it's my own look out!" These words were still on her lips, when she saw
        dowager lady Chia and the rest of the company arrive.}
}
{
}

\Verse{133}
{
    \zA{又不是个小孩儿，还这么淘气，仔细老太太说！”}
}
{
    \eA{They each sat where and how they pleased. First and foremost, a waiting-maid
        brought two trays of tea. After tea, lady Feng laid hold of a napkin,}
}
{
}

\Verse{134}
{
    \zA{鸳鸯笑 道：“很不与大奶奶相干，有我呢。”}
}
{
    \eA{made of foreign cloth, in which were wrapped a handful of blackwood
        chopsticks, encircled with three rings,}
}
{
}

\Verse{135}
{
    \zA{正说着，只见贾母等来了，各自随便坐下。}
}
{
    \eA{of inlaid silver, and distributed them on the tables, in the order in which
        they were placed.}
}
{
}

\Verse{136}
{
    \zA{先有丫鬟挨人递了茶。}
}
{
    \eA{"Bring that small hard-wood table over," old lady Chia then exclaimed;}
}
{
}

\Verse{137}
{
    \zA{大家吃毕， 凤姐手里拿着西洋布手巾，裹着一把乌木三镶银箸，按席罢下。}
}
{
    \eA{"and let our relative Liu sit next to me here!" No sooner did the servants
        hear her order than they hurried to move the table to where she wanted it.
        Lady Feng, during this interval,}
}
{
}

\Verse{138}
{
    \zA{贾母因说：“把那 一张小楠木桌子抬过来，让刘亲家挨着我这边坐。”}
}
{
    \eA{made a sign with her eye to Yüan Yang. Yüan Yang there and then dragged
        goody Liu out of the hall and began to impress in a low tone of voice
        various things on her mind.}
}
{
}

\Verse{139}
{
    \zA{众人听说，忙抬过来。}
}
{
    \eA{"This is the custom which prevails in our household,"}
}
{
}

\Verse{140}
{
    \zA{凤姐一 面递眼色与鸳鸯，鸳鸯便忙拉刘老老出去，悄悄的嘱咐了刘老老一席话，又说：“这 是我们家的规矩，要错了，我们就笑话呢。”}
}
{
    \eA{she proceeded, "and if you disregard it we'll have a laugh at your expense!"
        Having arranged everything she had in view, they at length returned to their
        places. Mrs. Hsüeh had come over, after her meal, so she simply seated
        herself on one side and sipped her tea. Dowager lady Chia with Pao-yü,}
}
{
}

\Verse{141}
{
    \zA{调停已毕，然后归坐。}
}
{
    \eA{Hsiang-yün, Tai-yü and Pao-ch'ai sat at one table.}
}
{
}

\Verse{142}
{
    \zA{薛姨妈是吃过 饭来的，不吃了，只坐在一边吃茶。}
}
{
    \eA{Madame Wang took the girls, Ying Ch'un, and her sisters, and occupied one
        table. Old goody Liu took a seat at a table next to dowager lady Chia.}
}
{
}

\Verse{143}
{
    \zA{贾母带着宝玉、湘云、黛玉、宝钗一桌，王夫 人带着迎春姐妹三人一桌，刘老老挨着贾母一桌。}
}
{
    \eA{Heretofore, while their old mistress had her repast, a young servant-maid
        usually stood by her to hold the finger bowl, yak-brush, napkin and other
        such necessaries, but Yüan Yang did not of late fulfil any of these duties,}
}
{
}

\Verse{144}
{
    \zA{贾母素日吃饭，皆有小丫鬟在旁 边拿着漱盂、麈尾、巾帕之物，如今鸳鸯是不当这差的了，今日偏接过麈尾来拂着。}
}
{
    \eA{so when, on this occasion, she deliberately seized the yak-brush and came
        over and flapped it about, the servant-girls concluded that she was bent
        upon playing some tricks upon goody Liu,}
}
{
}

\Verse{145}
{
    \zA{丫鬟们知他要捉弄刘老老，便躲开让他。}
}
{
    \eA{and they readily withdrew and let her have her way. While Yüan Yang attended
        to her self-imposed duties,}
}
{
}

\Verse{146}
{
    \zA{鸳鸯一面侍立，一面递眼色。}
}
{
    \eA{she winked at the old dame. "Miss," goody Liu exclaimed, "set your mind at
        ease!"}
}
{
}

\Verse{147}
{
    \zA{刘老老道： “姑娘放心。”}
}
{
    \eA{Goody Liu sat down at the table and took up the chopsticks,}
}
{
}

\Verse{148}
{
    \zA{那刘老老入了坐，拿起箸来，？}
}
{
    \eA{but so heavy and clumsy did she find them that she could not handle them
        conveniently.}
}
{
}

\Verse{149}
{
    \zA{甸甸的不伏手，原是凤姐和鸳鸯商议定了，单 拿了一双老年四楞象牙镶金的筷子给刘老老。}
}
{
    \eA{The fact is that lady Feng and Yüan Yang had put their heads together and
        decided to only assign to goody Liu a pair of antiquated four-cornered ivory
        chopsticks, inlaid with gold. "These forks," shouted goody Liu, after
        scrutinising them,}
}
{
}

\Verse{150}
{
    \zA{刘老老见了，说道：“这个叉巴子， 比我们那里的铁锨还？}
}
{
    \eA{"are heavier than the very iron-lever over at my place. How ever can I move
        them about?" This remark had the effect of making every one explode into a
        fit of laughter.}
}
{
}

\Verse{151}
{
    \zA{，那里拿的动他？”}
}
{
    \eA{But a married woman standing in the centre of the room,}
}
{
}

\Verse{152}
{
    \zA{说的众人都笑起来。}
}
{
    \eA{with a box in her hands, attracted their gaze. A waiting-maid went up to her
        and removed the cover of the box.}
}
{
}

\Verse{153}
{
    \zA{只见一个媳妇端了 一个盒子站在当地，一个丫鬟上来揭去盒盖，里面盛着两碗菜，李纨端了一碗放在 贾母桌上，凤姐偏拣了一碗鸽子蛋放在刘老老桌上。}
}
{
    \eA{Its contents were two bowls of eatables. Li Wan took one of these and placed
        it on dowager lady Chia's table, while lady Feng chose the bowl with
        pigeon's eggs and put it on goody Liu's table. "Please (commence)," Dowager
        lady Chia uttered from the near side, where she sat. Goody Liu at this
        speedily sprung to her feet.}
}
{
}

\Verse{154}
{
    \zA{贾母这边说声“请”，刘老老 便站起身来，高声说道：“老刘，老刘，食量大如牛。}
}
{
    \eA{"Old Liu, old Liu," she roared with a loud voice, "your eating capacity is
        as big as that of a buffalo! You've gorged like an old sow and can't raise
        your head up!" Then puffing out her cheeks,}
}
{
}

\Verse{155}
{
    \zA{吃个老母猪，不抬头！”}
}
{
    \eA{she added not a word. The whole party was at first taken quite aback.}
}
{
}

\Verse{156}
{
    \zA{说 完，却鼓着腮帮子，两眼直视，一声不语。}
}
{
    \eA{But, as soon as they heard the drift of her remarks, every one, both high as
        well as low, began to laugh boisterously.}
}
{
}

\Verse{157}
{
    \zA{众人先还发怔，后来一想，上上下下都 一齐哈哈大笑起来。}
}
{
    \eA{Hsiang-yün found it so difficult to restrain herself that she spurted out
        the tea she had in her mouth. Lin Tai-yü indulged in such laughter that she
        was quite out of breath,}
}
{
}

\Verse{158}
{
    \zA{湘云掌不住，一口茶都喷出来。}
}
{
    \eA{and propping herself up on the table, she kept on ejaculating 'Ai-yo.'}
}
{
}

\Verse{159}
{
    \zA{黛玉笑岔了气，伏着桌子只叫 “嗳哟”。}
}
{
    \eA{Pao-yü rolled into his grandmother's lap. The old lady herself was so amused
        that she clasped Pao-yü in her embrace,}
}
{
}

\Verse{160}
{
    \zA{宝玉滚到贾母怀里，贾母笑的搂着叫“心肝”。}
}
{
    \eA{and gave way to endearing epithets. Madame Wang laughed, and pointed at lady
        Feng with her finger; but as for saying a word,}
}
{
}

\Verse{161}
{
    \zA{王夫人笑的用手指着凤 姐儿，却说不出话来。}
}
{
    \eA{she could not. Mrs. Hsüeh had much difficulty in curbing her mirth, and she
        sputtered the tea, with which her mouth was full,}
}
{
}

\Verse{162}
{
    \zA{薛姨妈也掌不住，口里的茶喷了探春一裙子。}
}
{
    \eA{all over T'an Ch'un's petticoat. T'an Ch'un threw the contents of the
        teacup, she held in her hand, over Ying Ch'un; while Hsi Ch'un quitted her
        seat,}
}
{
}

\Verse{163}
{
    \zA{探春的茶碗都 合在迎春身上。}
}
{
    \eA{and, pulling her nurse away, bade her rub her stomach for her. Below, among
        the lower seats,}
}
{
}

\Verse{164}
{
    \zA{惜春离了坐位，拉着他奶母，叫“揉揉肠子”。}
}
{
    \eA{there was not one who was not with bent waist and doubled-up back. Some
        retired to a corner and, squatting down, laughed away.}
}
{
}
\ChapterImage{img/chapters/096第40回 劉姥姥初遊瀟湘館.jpg}


\Verse{165}
{
    \zA{地下无一个不弯腰 屈背，也有躲出去蹲着笑去的，也有忍着笑上来替他姐妹换衣裳的。}
}
{
    \eA{Others suppressed their laughter and came up and changed the clothes of
        their young mistresses. Lady Feng and Yuan Yang were the only ones, who kept
        their countenance. Still they continued helping old goody Liu to food.}
}
{
}

\Verse{166}
{
    \zA{独有凤姐鸳鸯 二人掌着，还只管让刘老老。}
}
{
    \eA{Old goody Liu took up the chopsticks. "Even the chickens in this place are
        fine," she went on to add, pretending,}
}
{
}

\Verse{167}
{
    \zA{刘老老拿起箸来，只觉不听使，又道：“这里的鸡儿也俊，下的这蛋也小巧， 怪俊的。}
}
{
    \eA{she did not hear what was going on; "the eggs they lay are small, but so
        dainty! How very pretty they are! Let me help myself to one!" The company
        had just managed to check themselves,}
}
{
}

\Verse{168}
{
    \zA{我且得一个儿！”}
}
{
    \eA{but, the moment these words fell on their ears,}
}
{
}

\Verse{169}
{
    \zA{众人方住了笑，听见这话，又笑起来。}
}
{
    \eA{they started again with their laughter. Old lady Chia laughed to such an
        extent that tears streamed from her eyes.}
}
{
}

\Verse{170}
{
    \zA{贾母笑的眼泪出 来只忍不住，琥珀在后捶着。}
}
{
    \eA{And so little could she bear the strain any longer that Hu Po stood behind
        her and patted her. "This must be the work of that vixen Feng!"}
}
{
}

\Verse{171}
{
    \zA{贾母笑道：“这定是凤丫头促狭鬼儿闹的!快别信他 的话了。”}
}
{
    \eA{old lady Chia laughed. "She has ever been up to tricks like a very imp, so
        be quick and disbelieve all her yarns!" Goody Liu was in the act of praising
        the eggs as small yet dainty,}
}
{
}

\Verse{172}
{
    \zA{那刘老老正夸鸡蛋小巧，凤姐儿笑道：“一两银子一个呢!你快尝尝罢， 冷了就不好吃了。”}
}
{
    \eA{when lady Feng interposed with a smile. "They're one tael each, be quick,
        and taste them;" she said; "they're not nice when they get cold!" Goody Liu
        forthwith stretched out the chopsticks with the intent of catching one;}
}
{
}

\Verse{173}
{
    \zA{刘老老便伸筷子要夹，那里夹的起来？}
}
{
    \eA{but how could she manage to do so? They rolled and rolled in the bowl for
        ever so long;}
}
{
}

\Verse{174}
{
    \zA{满碗里闹了一阵，好容 易撮起一个来，才伸着脖子要吃，偏又滑下来，滚在地下。}
}
{
    \eA{and, it was only after extreme difficulty that she succeeded in shoving one
        up. Extending her neck forward, she was about to put it in her mouth, when
        it slipped down again, and rolled on to the floor.}
}
{
}

\Verse{175}
{
    \zA{忙放下筷子要亲自去拣， 早有地下的人拣出去了。}
}
{
    \eA{She hastily banged down the chopsticks, and was going herself to pick it up,
        when a servant, who stood below, got hold of it and took it out of the room.}
}
{
}

\Verse{176}
{
    \zA{刘老老叹道：“一两银子，也没听见个响声儿就没了！”}
}
{
    \eA{Old goody Liu heaved a sigh. "A tael!" she soliloquised, "and here it goes
        without a sound!" Every one had long ago abandoned all idea of eating,}
}
{
}

\Verse{177}
{
    \zA{众人已没心吃饭，都看着他取笑。}
}
{
    \eA{and, gazing at her, they enjoyed the fun. "Who has now brought out these
        chopsticks again?"}
}
{
}

\Verse{178}
{
    \zA{贾母又说：“谁这会子又把那个筷子拿出来 了，又不请客摆大筵席!都是凤丫头支使的，还不换了呢。”}
}
{
    \eA{old lady Chia went on to ask. "We haven't invited any strangers or spread
        any large banquet! It must be that vixen Feng who gave them out! But don't
        you yet change them!" The servants, standing on the floor below, had indeed
        had no hand in getting those ivory chopsticks;}
}
{
}

\Verse{179}
{
    \zA{地下的人原不曾预备 这牙箸，本是凤姐和鸳鸯拿了来的，听如此说，忙收过去了，也照样换上一双乌木 镶银的。}
}
{
    \eA{they had, in fact, been brought by lady Feng and Yüan Yang; but when they
        heard these remarks, they hurried to put them away and to change them for a
        pair similar to those used by the others, made of blackwood inlaid with
        silver. "They've taken away the gold ones,"}
}
{
}

\Verse{180}
{
    \zA{刘老老道：“去了金的，又是银的，到底不及俺们那个伏手。”}
}
{
    \eA{old goody Liu shouted, "and here come silver ones! But, after all, they're
        not as handy as those we use!" "Should there be any poison in the viands,"}
}
{
}

\Verse{181}
{
    \zA{凤姐儿道： “菜里要有毒，这银子下去了就试的出来。”}
}
{
    \eA{lady Feng observed, "you can detect it, as soon as this silver is dipped
        into them!" "If there's poison in such viands as these,"}
}
{
}

\Verse{182}
{
    \zA{刘老老道：“这个菜里有毒，我们那 些都成了砒霜了!那怕毒死了，也要吃尽了。”}
}
{
    \eA{old goody Liu added, "why those of ours must be all arsenic! But though it
        be the death of me, I'll swallow every morsel!" Seeing how amusing the old
        woman was and with what relish she devoured her food,}
}
{
}

\Verse{183}
{
    \zA{贾母见他如此有趣，吃的又香甜， 把自己的菜也都端过来给他吃。}
}
{
    \eA{dowager lady Chia took her own dishes and passed them over to her. She then
        likewise bade an old matron take various viands and put them in a bowl for
        Pan Erh. But presently, the repast was concluded,}
}
{
}

\Verse{184}
{
    \zA{又命一个老嬷嬷来，将各样的菜给板儿夹在碗上。}
}
{
    \eA{and old lady Chia and all the other inmates adjoined into T'an Ch'un's
        bedroom for a chat. The remnants were, meanwhile, cleared away, and fresh
        tables were laid.}
}
{
}

\Verse{185}
{
    \zA{一时吃毕，贾母等都往探春卧室中去闲话，这里收拾残桌，又放了一桌。}
}
{
    \eA{Old goody Liu watched Li Wan and lady Feng sit opposite each other and eat.
        "Putting everything else aside," she sighed, "what most takes my fancy is
        the way things are done in your mansion.}
}
{
}

\Verse{186}
{
    \zA{刘老 老看着李纨与凤姐儿对坐着吃饭，叹道：“别的罢了，我只爱你们家这行事!怪道 说，‘礼出大家’。”}
}
{
    \eA{It isn't to be wondered at that the adage has it that: 'propriety originates
        from great families.'" "Don't be too touchy," lady Feng hastily smiled, "we
        all made fun of you just now." But barely had she done speaking,}
}
{
}

\Verse{187}
{
    \zA{凤姐儿忙笑道：“你可别多心，才刚不过大家取乐儿。”}
}
{
    \eA{when Yüan Yang too walked in. "Old goody Liu," she said laughingly, "don't
        be angry! I tender you my apologies, venerable dame!" "What are you saying,}
}
{
}

\Verse{188}
{
    \zA{一 言未了，鸳鸯也进来笑道：“老老别恼，我给你老人家赔个不是儿罢。”}
}
{
    \eA{Miss?" old goody Liu rejoined smiling. "We've coaxed our dowager lady to get
        a little distraction; and what reason is there to be angry? From the very
        first moment you spoke to me,}
}
{
}

\Verse{189}
{
    \zA{刘老老忙 笑道：“姑娘说那里的话？}
}
{
    \eA{I knew at once that it was intended to afford merriment to you all! Had I
        been angry at heart,}
}
{
}

\Verse{190}
{
    \zA{咱们哄着老太太开个心儿，有什么恼的!你先嘱咐我，我 就明白了，不过大家取笑儿。}
}
{
    \eA{I wouldn't have gone so far as to say what I did!" Yüan Yang then blew up
        the servants. "Why," she shouted, "don't you pour a cup of tea for the old
        dame?" "That sister-in-law," promptly explained old goody Liu,}
}
{
}

\Verse{191}
{
    \zA{我要恼，也就不说了。”}
}
{
    \eA{"gave me a cup a little while back. I've had it already. But you,}
}
{
}

\Verse{192}
{
    \zA{鸳鸯便骂人：“为什么不倒 茶给老老吃！”}
}
{
    \eA{Miss, must also have something to eat." Lady Feng dragged Yüan Yang into a
        seat. "Have your meal with us!"}
}
{
}

\Verse{193}
{
    \zA{刘老老忙道：“才刚那个嫂子倒了茶来，我吃过了，姑娘也该用饭 了。”}
}
{
    \eA{she said. "You'll thus save another fuss by and bye." Yüan Yang readily
        seated herself. The matrons came up and added to the number of bowls and
        chopsticks, and the trio went through their meal.}
}
{
}

\Verse{194}
{
    \zA{凤姐儿便拉鸳鸯坐下道：“你和我们吃罢，省了回来又闹。”}
}
{
    \eA{"From all I see," smiled goody Liu, "you people eat just a little and
        finish. It's lucky you don't feel the pangs of hunger! But it isn't
        astonishing if a whiff of wind can puff you over!"}
}
{
}

\Verse{195}
{
    \zA{鸳鸯便坐下了， 婆子们添上碗箸来，三人吃毕。}
}
{
    \eA{"A good many eatables remained over to-day. Where are they all gone to?"
        Yüan Yang inquired. "They haven't as yet been apportioned!"}
}
{
}

\Verse{196}
{
    \zA{刘老老笑道：“我看你们这些人，都只吃这一点儿 就完了，亏你们也不饿。}
}
{
    \eA{the matrons responded. "They're kept in here until they can be given in a
        lump to them to eat!" "They can't get through so many things!" Yüan Yang
        resumed. "You had as well therefore choose two bowls and send them over to
        that girl P'ing,}
}
{
}

\Verse{197}
{
    \zA{怪道风儿都吹的倒！”}
}
{
    \eA{in your mistress Secundus' rooms." "She has had her repast long ago."}
}
{
}

\Verse{198}
{
    \zA{鸳鸯便问：“今儿剩的不少，都那 里去了？”}
}
{
    \eA{lady Feng put in. "There's no need to give her any!" "With what she can't
        eat, herself," Yüan Yang continued,}
}
{
}

\Verse{199}
{
    \zA{婆子们道：“都还没散呢，在这里等着，一齐散给他们吃。”}
}
{
    \eA{"she can feed the cats." At these words, a matron lost no time in selecting
        two sorts of eatables, and, taking the box, she went to take them over.}
}
{
}

\Verse{200}
{
    \zA{鸳鸯道： “他们吃不了这些，挑两碗给二奶奶屋里平丫头送去。”}
}
{
    \eA{"Where's Su Yun gone to?" Yüan Yang asked. "They're all in here having their
        meal together." Li Wan replied. "What do you want her for again?" "Well, in
        that case, never mind,"}
}
{
}

\Verse{201}
{
    \zA{凤姐道：“他早吃了饭了， 不用给他。”}
}
{
    \eA{Yüan Yang answered. "Hsi Jen isn't here," lady Feng observed, "so tell some
        one to take her a few things!"}
}
{
}

\Verse{202}
{
    \zA{鸳鸯道：“他吃不了，喂你的猫。”}
}
{
    \eA{Yuan Yang, hearing this, directed a servant to send her also a few eatables.}
}
{
}

\Verse{203}
{
    \zA{婆子听了，忙拣了两样，拿盒子 送去。}
}
{
    \eA{"Have the partition boxes been filled with wine for by and bye?" Yüan Yang
        went on to ask the matrons.}
}
{
}

\Verse{204}
{
    \zA{鸳鸯道：“素云那里去了？”}
}
{
    \eA{"They'll be ready, I think, in a little while," a matron explained.}
}
{
}

\Verse{205}
{
    \zA{李纨道：“他们都在这里一处吃，又找他做什 么？”}
}
{
    \eA{"Hurry them up a bit!" Yüan Yang added. The matron signified her assent.
        Lady Feng and her friends then came into T'an Ch'un's apartments,}
}
{
}

\Verse{206}
{
    \zA{鸳鸯道：“这就罢了。”}
}
{
    \eA{where they found the ladies chatting and laughing.}
}
{
}

\Verse{207}
{
    \zA{凤姐道：“袭人不在这里，你倒是叫人送两样给他 去。”}
}
{
    \eA{T'an Ch'un had ever shown an inclination for plenty of room. Hence that
        suite of three apartments had never been partitioned. In the centre was
        placed a large table of rosewood and Ta li marble.}
}
{
}

\Verse{208}
{
    \zA{鸳鸯听说，便命人也送两样去。}
}
{
    \eA{On this table, were laid in a heap every kind of copyslips written by
        persons of note.}
}
{
}

\Verse{209}
{
    \zA{鸳鸯又问婆子们：“回来吃酒的攒盒，可装 上了？”}
}
{
    \eA{Several tens of valuable inkslabs and various specimens of tubes and
        receptacles for pens figured also about; the pens in which were as thickly
        packed as trees in a forest.}
}
{
}

\Verse{210}
{
    \zA{婆子道：“想必还得一会子。”}
}
{
    \eA{On the off side, stood a flower bowl from the 'Ju' kiln, as large as a
        bushel measure.}
}
{
}

\Verse{211}
{
    \zA{鸳鸯道：“催着些儿。”}
}
{
    \eA{In it was placed, till it was quite full, a bunch of white chrysanthemums,}
}
{
}

\Verse{212}
{
    \zA{婆子答应了。}
}
{
    \eA{in appearance like crystal balls.}
}
{
}

\Verse{213}
{
    \zA{凤姐等来至探春房中，只见他娘儿们正说笑。}
}
{
    \eA{In the middle of the west wall, was suspended a large picture representing
        vapor and rain; the handiwork of Mi Nang-yang.}
}
{
}

\Verse{214}
{
    \zA{探春素喜阔朗，这三间屋子并不 曾隔断，当地放着一张花梨大理石大案，案上堆着各种名人法帖，并数十方宝砚， 各色笔筒，笔海内插的笔如树林一般。}
}
{
    \eA{On the left and right of this picture was hung a pair of antithetical
        scrolls—the autograph of Yen Lü. The lines on these scrolls were: Wild
        scenes are to the taste of those who leisure love, And springs and rookeries
        are their rustic resort. On the table, figured a large tripod. On the left,
        stood on a blackwood cabinet,}
}
{
}

\Verse{215}
{
    \zA{那一边设着斗大的一个汝窑花囊，插着满满 的一囊水晶球的白菊。}
}
{
    \eA{a huge bowl from a renowned government kiln. This bowl contained about ten
        "Buddha's hands" of beautiful yellow and fine proportions. On the right, was
        suspended,}
}
{
}

\Verse{216}
{
    \zA{西墙上当中挂着一大幅米襄阳《烟雨图》。}
}
{
    \eA{on a Japanese-lacquered frame, a white jade sonorous plate. Its shape
        resembled two eyes, one by the side of the other.}
}
{
}

\Verse{217}
{
    \zA{左右挂着一副对 联，乃是颜鲁公墨迹。}
}
{
    \eA{Next to it hung a small hammer. Pan Erh had become a little more confident
        and was about to seize the hammer and beat the plate,}
}
{
}

\Verse{218}
{
    \zA{其联云： 烟霞闲骨格， 泉石野生涯。}
}
{
    \eA{when the waiting-maids hastened to prevent him. Next, he wanted a "Buddha's
        hand" to eat.}
}
{
}

\Verse{219}
{
    \zA{案上设着大鼎，左边紫檀架上放着一个大官窑的大盘，盘内盛着数十个娇黄玲珑大 佛手。}
}
{
    \eA{T'an Ch'un chose one and let him have it. "You may play with it," she said,
        "but you can't eat it." On the east side stood a sleeping divan. On a
        movable bed was hung a leek-green gauze curtain, ornamented with double
        embroideries,}
}
{
}

\Verse{220}
{
    \zA{右边洋漆架上悬着一个白玉比目磬，傍边挂着小槌。}
}
{
    \eA{representing flowers, plants and insects. Pan Erh ran up to have a look.
        "This is a green-cicada," he shouted; "this a grasshopper!"}
}
{
}

\Verse{221}
{
    \zA{那板儿略熟了些，便要 摘那槌子去击，丫鬟们忙拦住他。}
}
{
    \eA{But old goody Liu promptly gave him a slap. "You mean scamp!" she cried.
        "What an awful rumpus you're kicking up! I simply brought you along with me
        to look at things;}
}
{
}

\Verse{222}
{
    \zA{他又要那佛手吃，探春拣了一个给他，说：“玩 罢，吃不得的。”}
}
{
    \eA{and lo, you put on airs;" and she beat Pan Erh until he burst out crying. It
        was only after every one quickly combined in using their efforts to solace
        him that he at length desisted.}
}
{
}

\Verse{223}
{
    \zA{东边便设着卧榻拔步床，上悬着葱绿双绣花卉草虫的纱帐。}
}
{
    \eA{Old lady Chia then looked through the gauze casement into the back court for
        some time. "The dryandra trees by the eaves of the covered passage are
        growing all right,"}
}
{
}

\Verse{224}
{
    \zA{板儿 又跑来看，说：“这是蝈蝈，这是蚂蚱。”}
}
{
    \eA{she remarked. "The only thing is that their foliage is rather sparse." But
        while she passed this remark, a sudden gust of wind swept by,}
}
{
}

\Verse{225}
{
    \zA{刘老老忙打了他一巴掌，道：“下作黄 子!没干没净的乱闹。}
}
{
    \eA{and faintly on her ear fell the strains of music. "In whose house is there a
        wedding?" old lady Chia inquired. "This place must be very near the street!"}
}
{
}

\Verse{226}
{
    \zA{倒叫你进来瞧瞧，就上脸了！”}
}
{
    \eA{"How could one hear what's going on in the street?" Madame Wang and the
        others smiled.}
}
{
}

\Verse{227}
{
    \zA{打的板儿哭起来，众人忙劝 解方罢。}
}
{
    \eA{"It's our twelve girls practising on their wind and string instruments!" "As
        they're practising,"}
}
{
}

\Verse{228}
{
    \zA{贾母隔着纱窗后往院内看了一回，因说道：“后廊檐下的梧桐也好了，只是细 些。”}
}
{
    \eA{dowager lady Chia eagerly cried, smilingly, "why not ask them to come in
        here and practise? They'll be able to have a stroll also, while we, on our
        part, will derive some enjoyment."}
}
{
}

\Verse{229}
{
    \zA{正说话，忽一阵风过，隐隐听得鼓乐之声。}
}
{
    \eA{Upon hearing this suggestion, lady Feng immediately directed a servant to go
        out and call them in.}
}
{
}

\Verse{230}
{
    \zA{贾母问：“是谁家娶亲呢？}
}
{
    \eA{She further issued orders to bring a table and spread a red cover over it.}
}
{
}

\Verse{231}
{
    \zA{这里 临街倒近。”}
}
{
    \eA{"Let it be put," old lady Chia chimed in,}
}
{
}

\Verse{232}
{
    \zA{王夫人等笑回道：“街上的那里听的见？}
}
{
    \eA{"in the water-pavilion of the Lotus Fragrance Arbour, for (the music) will
        borrow the ripple of the stream and sound ever so much more pleasant to the
        ear.}
}
{
}

\Verse{233}
{
    \zA{这是咱们的那十来个女孩子 们演习吹打呢。”}
}
{
    \eA{We can by and bye drink our wine in the Cho Chin Hall; we'll thus have ample
        room, and be able to listen from close!"}
}
{
}

\Verse{234}
{
    \zA{贾母便笑道：“既他们演，何不叫他们进来演习，他们也逛一逛， 咱们也乐了，不好吗？”}
}
{
    \eA{Every one admitted that the spot was well adapted. Dowager lady Chia turned
        herself towards Mrs. Hsüeh. "Let's get ahead!" she laughed.}
}
{
}

\Verse{235}
{
    \zA{凤姐听说，忙命人出去叫来，赶着吩咐摆下条桌，铺上红 毡子。}
}
{
    \eA{"The young ladies don't like any one to come in here, for fear lest their
        quarters should get contaminated; so don't let us show ourselves
        disregardful of their wishes!}
}
{
}

\Verse{236}
{
    \zA{贾母道：“就铺排在藕香榭的水亭子上，借着水音更好听。}
}
{
    \eA{The right thing would be to go and have our wine aboard one of those boats!"
        As she spoke, one and all rose to their feet.}
}
{
}

\Verse{237}
{
    \zA{回来咱们就在缀 锦阁底下吃酒，又宽阔，又听的近。”}
}
{
    \eA{They were making their way out when T'an Ch'un interposed. "What's this that
        you're saying?" she smiled. "Please do seat yourselves,}
}
{
}

\Verse{238}
{
    \zA{众人都说好。}
}
{
    \eA{venerable senior, and you,}
}
{
}

\Verse{239}
{
    \zA{贾母向薛姨妈笑道：“咱们走 罢，他们姐妹们都不大喜欢人来，生怕腌？}
}
{
    \eA{Mrs. Hsüeh, and Madame Wang! You can't be going yet?" "These three girls of
        mine are really nice! There are only two mistresses that are simply
        dreadful." Dowager lady Chia said smilingly.}
}
{
}

\Verse{240}
{
    \zA{了屋子。}
}
{
    \eA{"When we get drunk shortly,}
}
{
}

\Verse{241}
{
    \zA{咱们别没眼色儿，正经坐会子 船，喝酒去罢。”}
}
{
    \eA{we'll go and sit in their rooms and have a lark!" These words evoked
        laughter from every one. In a body they quitted the place.}
}
{
}

\Verse{242}
{
    \zA{说着，大家起身便走。}
}
{
    \eA{But they had not proceeded far before they reached the bank covered with
        aquatic plants,}
}
{
}

\Verse{243}
{
    \zA{探春笑道：“这是那里的话？}
}
{
    \eA{to which place the boat-women, who had been brought from Ku Su,}
}
{
}

\Verse{244}
{
    \zA{求着老太太、 姨妈、太太来坐坐还不能呢！”}
}
{
    \eA{had already punted two crab-wood boats. Into one of these boats, they helped
        old lady Chia, Madame Wang, Mrs. Hsüeh,}
}
{
}

\Verse{245}
{
    \zA{贾母笑道：“我的这三丫头倒好，只有两个玉儿可 恶。}
}
{
    \eA{old goody Liu, Yüan Yang, and Yü Ch'uan-Erh. Last in order Li Wan followed
        on board. But lady Feng too stepped in, and standing up on the bow,}
}
{
}

\Verse{246}
{
    \zA{回来喝醉了，咱们偏往他们屋里闹去！”}
}
{
    \eA{she insisted upon punting. Dowager lady Chia, however, remonstrated from her
        seat in the bottom of the boat.}
}
{
}

\Verse{247}
{
    \zA{说着众人都笑了。}
}
{
    \eA{"This isn't a joke," she cried, "we're not on the river,}
}
{
}
\ChapterImage{img/chapters/097第40回 劉姥姥初遊紫菱洲.jpg}


\Verse{248}
{
    \zA{一齐出来走不多远，已到了荇叶渚，那姑苏选来的几个驾娘早把两只棠木舫撑 来。}
}
{
    \eA{it's true, but there are some very deep places about, so be quick and come
        in. Do it for my sake." "What's there to be afraid of?" lady Feng laughed.
        "Compose your mind,}
}
{
}

\Verse{249}
{
    \zA{众人扶了贾母，王夫人、薛姨妈、刘老老、鸳鸯、玉钏儿上了这一只船，次后 李纨也跟上去。}
}
{
    \eA{worthy ancestor." Saying this, the boat was pushed off with one shove. When
        it reached the middle of the lake, lady Feng became nervous, for the craft
        was small and the occupants many,}
}
{
}

\Verse{250}
{
    \zA{凤姐也上去，立在船头上，也要撑船。}
}
{
    \eA{and hastily handing the pole to a boatwoman, she squatted down at last. Ying
        Ch'un, her sisters,}
}
{
}

\Verse{251}
{
    \zA{贾母在舱内道：“那不是玩 的!虽不是河里，也有好深的，你快给我进来。”}
}
{
    \eA{their cousins, as well as Pao-yü subsequently got on board the second boat,
        and followed in their track; while the rest of the company, consisting of
        old nurses and a bevy of waiting-maids,}
}
{
}

\Verse{252}
{
    \zA{凤姐笑道：“怕什么!老祖宗只管 放心。”}
}
{
    \eA{kept pace with them along the bank of the stream. "All these broken lotus
        leaves are dreadful!"}
}
{
}

\Verse{253}
{
    \zA{说着，便一篙点开，到了池当中。}
}
{
    \eA{Pao-yü shouted. "Why don't you yet tell the servants to pull them off?"}
}
{
}

\Verse{254}
{
    \zA{船小人多，凤姐只觉乱晃，忙把篙子递 与驾娘，方蹲下去。}
}
{
    \eA{"When was this garden left quiet during all the days of this year?"
        Pao-ch'ai smiled. "Why, people have come, day after day, to visit it,}
}
{
}

\Verse{255}
{
    \zA{然后迎春姐妹等并宝玉上了那只，随后跟来。}
}
{
    \eA{so was there ever any time to tell the servants to come and clean it?" "I
        have the greatest abhorrence," Lin Tai-yü chimed in, "for Li I's poetical
        works,}
}
{
}

\Verse{256}
{
    \zA{其馀老嬷嬷众丫 鬟俱沿河随行。}
}
{
    \eA{but there's only this line in them which I like: "'Leave the dry lotus
        leaves so as to hear the patter of the rain.'}
}
{
}

\Verse{257}
{
    \zA{宝玉道：“这些破荷叶可恨，怎么还不叫人来拔去？”}
}
{
    \eA{"and here you people deliberately mean again not to leave the dry lotus stay
        where they are." "This is indeed a fine line!"}
}
{
}

\Verse{258}
{
    \zA{宝钗笑道： “今年这几日，何曾饶了这园子闲了一闲，天天逛，那里还有叫人来收拾的工夫 呢？”}
}
{
    \eA{Pao-yü exclaimed. "We mustn't hereafter let them pull them away!" While this
        conversation continued, they reached the shoaly inlet under the flower-laden
        beech. They felt a coolness from the shady overgrowth penetrate their very
        bones.}
}
{
}

\Verse{259}
{
    \zA{黛玉道：“我最不喜欢李义山的诗，只喜他这一句：‘留得残荷听雨声。’}
}
{
    \eA{The decaying vegetation and the withered aquatic chestnut plants on the
        sand-bank enhanced, to a greater degree, the beauty of the autumn scenery.}
}
{
}

\Verse{260}
{
    \zA{偏你们又不留着残荷了。”}
}
{
    \eA{Dowager lady Chia at this point observed some spotless rooms on the bank,}
}
{
}

\Verse{261}
{
    \zA{宝玉道：“果然好句，以后咱们别叫拔去了。”}
}
{
    \eA{so spick and so span. "Are not these Miss Hsüeh's quarters," she asked.
        "Eh?" "Yes, they are!"}
}
{
}

\Verse{262}
{
    \zA{说着已到了花溆的萝港之下，觉得阴森透骨，两滩上衰草残菱，更助秋兴。}
}
{
    \eA{everybody answered. Old lady Chia promptly bade them go alongside, and
        wending their way up the marble steps, which seemed to lead to the clouds,
        they in a body entered the Heng Wu court.}
}
{
}

\Verse{263}
{
    \zA{贾 母因见岸上的清厦旷朗，便问：“这是薛姑娘的屋子不是？”}
}
{
    \eA{Here they felt a peculiar perfume come wafting into their nostrils, for the
        colder the season got the greener grew that strange vegetation, and those
        fairy-like creepers.}
}
{
}

\Verse{264}
{
    \zA{众人道：“是。”}
}
{
    \eA{The various plants were laden with seeds,}
}
{
}

\Verse{265}
{
    \zA{贾 母忙命拢岸，顺着云步石梯上去，一同进了蘅芜院。}
}
{
    \eA{which closely resembled red coral beans, as they drooped in lovely clusters.
        The house, as soon as they put their foot into it,}
}
{
}

\Verse{266}
{
    \zA{只觉异香扑鼻，那些奇草仙藤， 愈冷愈苍翠，都结了实，似珊瑚豆子一般，累垂可爱。}
}
{
    \eA{presented the aspect of a snow cave. There was a total absence of every
        object of ornament. On the table figured merely an earthenware vase, in
        which were placed several chrysanthemums.}
}
{
}

\Verse{267}
{
    \zA{及进了房屋，雪洞一般，一 色的玩器全无。}
}
{
    \eA{A few books and teacups were also conspicuous, but no further knicknacks. On
        the bed was suspended a green gauze curtain,}
}
{
}

\Verse{268}
{
    \zA{案上止有一个土定瓶，瓶中供着数枝菊，并两部书，茶奁、茶杯而 已。}
}
{
    \eA{and of equally extreme plainness were the coverlets and mattresses belonging
        to it. "This child," dowager lady Chia sighed, "is too simple! If you've got
        nothing to lay about,}
}
{
}

\Verse{269}
{
    \zA{床上只吊着青纱帐幔，衾褥也十分朴素。}
}
{
    \eA{why not ask your aunt for a few articles? I would never raise any objection.
        I never thought about them.}
}
{
}

\Verse{270}
{
    \zA{贾母叹道：“这孩子太老实了!你没 有陈设，何妨和你姨娘要些？}
}
{
    \eA{Your things, of course, have been left at home, and have not been brought
        over." So saying, she told Yuan Yang to go and fetch several bric-a-brac.
        She next went on to call lady Feng to task.}
}
{
}

\Verse{271}
{
    \zA{我也没理论，也没想到。}
}
{
    \eA{"She herself wouldn't have them," (lady Feng) rejoined. "We really sent over
        a few,}
}
{
}

\Verse{272}
{
    \zA{你们的东西，自然在家里没 带了来。”}
}
{
    \eA{but she refused every one of them and returned them." "In her home also,"
        smiled Mrs. Hsüeh,}
}
{
}

\Verse{273}
{
    \zA{说着，命鸳鸯去取些古董来，又嗔着凤姐儿：“不送些玩器来给你妹妹， 这样小器！”}
}
{
    \eA{"she does not go in very much for such sort of things." Old lady Chia nodded
        her head. "It will never do!" she added. "It does, it's true, save trouble;
        but were some relative to come on a visit,}
}
{
}

\Verse{274}
{
    \zA{王夫人凤姐等都笑回说：“他自己不要么，我们原送了来，都退回去 了。”}
}
{
    \eA{she'll find things in an impossible way. In the second place, such
        simplicity in the apartments of young ladies of tender age is quite
        unpropitious! Why, if you young people go on in this way,}
}
{
}

\Verse{275}
{
    \zA{薛姨妈也笑说道：“他在家里也不大弄这些东西。”}
}
{
    \eA{we old fogies should go further and live in stables! You've all heard what
        is said in those books and plays about the dreadful luxury,}
}
{
}

\Verse{276}
{
    \zA{贾母摇头道：“那使不 得。}
}
{
    \eA{with which young ladies' quarters are got up.}
}
{
}

\Verse{277}
{
    \zA{虽然他省事，倘或来个亲戚，看着不像，二则年轻的姑娘们，屋里这么素净， 也忌讳。}
}
{
    \eA{And though these girls of ours could not presume to place themselves on the
        same footing as those young ladies, they shouldn't nevertheless exceed too
        much the bounds of what constitutes the right thing. If they have any
        objects ready at hand, why shouldn't they lay them out?}
}
{
}

\Verse{278}
{
    \zA{我们这老婆子，越发该住马圈去了。}
}
{
    \eA{And if they have any strong predilection for simplicity, a few things less
        will do quite as well.}
}
{
}

\Verse{279}
{
    \zA{你们听那些书上戏上说的小姐们的绣 房，精致的还了得呢!他们姐妹们虽不敢比那些小姐们，也别很离了格儿。}
}
{
    \eA{I've always had the greatest knack for titifying a room, but being an old
        woman now I haven't the ease and inclination to attend to such things! These
        girls are, however, learning how to do things very nicely. I was afraid that
        there would be an appearance of vulgarity in what they did,}
}
{
}

\Verse{280}
{
    \zA{有现成 的东西，为什么不摆呢？}
}
{
    \eA{and that, even had they anything worth having, they'd so place them about as
        to spoil them;}
}
{
}

\Verse{281}
{
    \zA{要很爱素净，少几样倒使得。}
}
{
    \eA{but from what I can see there's nothing vulgar about them. But let me now
        put things right for you,}
}
{
}

\Verse{282}
{
    \zA{我最会收拾屋子，如今老了， 没这个闲心了。}
}
{
    \eA{and I'll wager that everything will look grand as well as plain. I've got a
        couple of my own knicknacks, which I've managed to keep to this day,}
}
{
}

\Verse{283}
{
    \zA{他们姐妹们也还学着收拾的好。}
}
{
    \eA{by not allowing Pao-yü to get a glimpse of them; for had he ever seen them,}
}
{
}

\Verse{284}
{
    \zA{只怕俗气，有好东西也摆坏了。}
}
{
    \eA{they too would have long ago disappeared!" Continuing, she called Yüan Yang.}
}
{
}

\Verse{285}
{
    \zA{我 看他们还不俗。}
}
{
    \eA{"Fetch that marble pot with scenery on it,"}
}
{
}

\Verse{286}
{
    \zA{如今等我替你收拾，包管又大方又素净。}
}
{
    \eA{she said to her; "that gauze screen, and that tripod of transparent stone
        with black streaks, which you'll find in there,}
}
{
}

\Verse{287}
{
    \zA{我的两件体己，收到如今， 没给宝玉看见过，若经了他的眼也没了。”}
}
{
    \eA{and lay out all three on this table. They'll be ample! Bring likewise those
        ink pictures and white silk curtains, and change these curtains." Yüan Yang
        expressed her obedience.}
}
{
}

\Verse{288}
{
    \zA{说着，叫过鸳鸯来，吩咐道：“你把那 石头盆景儿和那架纱照屏，还有个墨烟冻石鼎拿来：这三样摆在这案上就够了。}
}
{
    \eA{"All these articles have been put away in the eastern loft," she smiled. "In
        what boxes they've been put, I couldn't tell; I must therefore go and find
        them quietly and if I bring them over to-morrow, it will be time enough."}
}
{
}

\Verse{289}
{
    \zA{再 把那水墨字画白绫帐子拿来，把这帐子也换了。”}
}
{
    \eA{"To-morrow or the day after will do very well; but don't forget, that's
        all," dowager lady Chia urged. While conversing, they sat for a while.}
}
{
}

\Verse{290}
{
    \zA{鸳鸯答应着，笑道：“这些东西 都搁在东楼上不知那个箱子里，还得慢慢找去，明儿再拿去也罢了。”}
}
{
    \eA{Presently, they left the rooms and repaired straightway into the Cho Chin
        hall. Wen Kuan and the other girls came up and paid their obeisance. They
        next inquired what songs they were to practise. "You'd better choose a few
        pieces to rehearse out of those you know best,"}
}
{
}

\Verse{291}
{
    \zA{贾母道：“明 日后日都使得，只别忘了。”}
}
{
    \eA{old lady Chia rejoined. Wen Kuan and her companions then withdrew and betook
        themselves to the Lotus Fragrance Pavilion.}
}
{
}

\Verse{292}
{
    \zA{说着，坐了一回，方出来，一径来至缀锦阁下。}
}
{
    \eA{But we will leave them there without further allusion to them. During this
        while, lady Feng had already, with the help of servants,}
}
{
}

\Verse{293}
{
    \zA{文官等上来请过安，因问：“演 习何曲？”}
}
{
    \eA{got everything in perfect order. On the left and right of the side of honour
        were placed two divans.}
}
{
}

\Verse{294}
{
    \zA{贾母道：“只拣你们熟的演习几套罢。”}
}
{
    \eA{These divans were completely covered with embroidered covers and fine
        variegated mats.}
}
{
}

\Verse{295}
{
    \zA{文官等下来，往藕香榭去不提。}
}
{
    \eA{In front of each divan stood two lacquer teapoys, inlaid, some with designs
        of crab-apple flowers;}
}
{
}

\Verse{296}
{
    \zA{这里凤姐已带着人摆设齐整，上面左右两张榻，榻上都铺着锦？}
}
{
    \eA{others of plum blossom, some of lotus leaves, others of sun-flowers. Some of
        these teapoys were square, others round. Their shapes were all different.}
}
{
}

\Verse{297}
{
    \zA{蓉簟，每一榻前两 张雕漆几，也有海棠式的，也有梅花式的，也有荷叶式的，也有葵花式的，也有方 的，有圆的，其式不一。}
}
{
    \eA{On each was placed a set consisting of a stove and a bottle, also a box with
        partitions. The two divans and four teapoys, in the place of honour, were
        used by dowager lady Chia and Mrs. Hsüeh. The chair and two teapoys in the
        next best place,}
}
{
}

\Verse{298}
{
    \zA{一个上头放着一分炉瓶，一个攒盒。}
}
{
    \eA{by Madame Wang. The rest of the inmates had, all alike, a chair and a
        teapoy.}
}
{
}

\Verse{299}
{
    \zA{上面二榻四几，是贾 母薛姨妈；下面一椅两几，是王夫人的。}
}
{
    \eA{On the east side sat old goody Liu. Below old goody Liu came Madame Wang. On
        the west was seated Shih Hsiang-yün. The second place was occupied by
        Pao-ch'ai;}
}
{
}

\Verse{300}
{
    \zA{馀者都是一椅一几。}
}
{
    \eA{the third by Tai-yü; the fourth by Ying Ch'un.}
}
{
}

\Verse{301}
{
    \zA{东边刘老老，刘老老 之下便是王夫人。}
}
{
    \eA{T'an Ch'un and Hsi Ch'un filled the lower seats, in their proper order;
        Pao-yü sat in the last place.}
}
{
}

\Verse{302}
{
    \zA{西边便是湘云，第二便是宝钗，第三便是黛玉，第四迎春，探春 惜春挨次排下去，宝玉在末。}
}
{
    \eA{The two teapoys assigned to Li Wan and lady Feng stood within the third line
        of railings, and beyond the second row of gauze frames. The pattern of the
        partition-boxes corresponded likewise with the pattern on the teapoys.}
}
{
}

\Verse{303}
{
    \zA{李纨凤姐二人之几设于三层槛内、二层纱厨之外。}
}
{
    \eA{Each inmate had a black decanter, with silver, inlaid in foreign designs; as
        well as an ornamented,}
}
{
}

\Verse{304}
{
    \zA{攒 盒式样，亦随几之式样。}
}
{
    \eA{enamelled cup. After they had all occupied the seats assigned to them,}
}
{
}

\Verse{305}
{
    \zA{每人一把乌银洋錾自斟壶，一个十锦珐琅杯。}
}
{
    \eA{dowager lady Chia took the initiative and smilingly suggested: "Let's begin
        by drinking a couple of cups of wine.}
}
{
}

\Verse{306}
{
    \zA{大家坐定，贾母先笑道：“咱们先吃两杯，今日也行一个令，才有意思。”}
}
{
    \eA{But we should also have a game of forfeits to-day, we'll have plenty of fun
        then." "You, venerable senior, must certainly have a good wine order to
        impose,"}
}
{
}

\Verse{307}
{
    \zA{薛 姨妈笑说道：“老太太自然有好酒令，我们如何会呢!安心叫我们醉了。}
}
{
    \eA{Mrs. Hsüeh laughingly observed, "but how could we ever comply with it? But
        if your aim be to intoxicate us, why, we'll all straightway drink one or two
        cups more than is good for us and finish!"}
}
{
}

\Verse{308}
{
    \zA{我们都多 吃两杯就有了。”}
}
{
    \eA{"Here's Mrs. Hsüeh beginning to be modest again to-day!"}
}
{
}

\Verse{309}
{
    \zA{贾母笑道：“姨太太今儿也过谦起来，想是厌我老了。”}
}
{
    \eA{old lady Chia smiled. "But I expect it's because she looks down upon me as
        being an old hag!" "It isn't modesty!" Mrs. Hsüeh replied smiling.}
}
{
}

\Verse{310}
{
    \zA{薛姨妈 笑道：“不是谦，只怕行不上来，倒是笑话了。”}
}
{
    \eA{"It's all a dread lest I shouldn't be able to observe the order and thus
        incur ridicule." "If you don't give the right answer,"}
}
{
}

\Verse{311}
{
    \zA{王夫人忙笑道：“便说不上来， 只多吃了一杯酒，醉了睡觉去，还有谁笑话咱们不成。”}
}
{
    \eA{Madame Wang promptly interposed with a smile, "you'll only have to drink a
        cup or two more of wine, and should we get drunk, we can go to sleep; and
        who'll, pray laugh at us?" Mrs. Hsüeh nodded her head.}
}
{
}

\Verse{312}
{
    \zA{薛姨妈点头笑道：“依令。}
}
{
    \eA{"I'll agree to the order," she laughed, "but, dear senior,}
}
{
}

\Verse{313}
{
    \zA{老太太到底吃一杯令酒才是。”}
}
{
    \eA{you must, after all, do the right thing and have a cup of wine to start it."}
}
{
}

\Verse{314}
{
    \zA{贾母笑道：“这个自然。”}
}
{
    \eA{"This is quite natural!" old lady Chia answered laughingly;}
}
{
}

\Verse{315}
{
    \zA{说着便吃了一杯。}
}
{
    \eA{and with these words, she forthwith emptied a cup.}
}
{
}

\Verse{316}
{
    \zA{凤姐 儿忙走至当地，笑道：“既行令，还叫鸳鸯姐姐来行才好。”}
}
{
    \eA{Lady Feng with hurried steps advanced to the centre of the room. "If we are
        to play at forfeits," she smilingly proposed, "we'd better invite sister
        Yüan Yang to come and join us."}
}
{
}

\Verse{317}
{
    \zA{众人都知贾母所行之 令，必得鸳鸯提着，故听了这话都说很是。}
}
{
    \eA{The whole company was perfectly aware that if dowager lady Chia had to give
        out the rule of forfeits, Yüan Yang would necessarily have to suggest it, so
        the moment they heard the proposal they,}
}
{
}

\Verse{318}
{
    \zA{凤姐便拉着鸳鸯过来。}
}
{
    \eA{with common consent, approved it as excellent.}
}
{
}

\Verse{319}
{
    \zA{王夫人笑道：“既 在令内，没有站着的理。”}
}
{
    \eA{Lady Feng therefore there and then dragged Yüan Yang over. "As you're to
        take a part in the game of forfeits,"}
}
{
}

\Verse{320}
{
    \zA{回头命小丫头子：“端一张椅子，放在你二位奶奶的席 上。”}
}
{
    \eA{Madame Wang smilingly observed, "there's no reason why you should stand up."
        And turning her head round, "Bring over," she bade a young waiting-maid,}
}
{
}

\Verse{321}
{
    \zA{鸳鸯也半推半就，谢了坐便坐下，也吃了一钟酒，笑道：“酒令大如军令。}
}
{
    \eA{"a chair and place it at your Mistress Secunda's table." Yüan Yang, half
        refusing and half assenting, expressed her thanks, and took the seat. After
        partaking also of a cup of wine,}
}
{
}

\Verse{322}
{
    \zA{不论尊卑，惟我是主，违了我的话，是要受罚的。”}
}
{
    \eA{"Drinking rules," she smiled, "resemble very much martial law; so
        irrespective of high or low, I alone will preside. Any one therefore who
        disobeys my words will have to suffer a penalty."}
}
{
}

\Verse{323}
{
    \zA{王夫人等都笑道：“一定如此， 快些说。”}
}
{
    \eA{"Of course, it should be so!" Madame Wang and the others laughed, "so be
        quick and give out the rule!"}
}
{
}

\Verse{324}
{
    \zA{鸳鸯未开口，刘老老便下席，摆手道：“别这样捉弄人!我家去了。”}
}
{
    \eA{But before Yüan Yang had as yet opened her lips to speak, old goody Liu left
        the table, and waving her hand: "Don't," she said, "make fun of people in
        this way,}
}
{
}

\Verse{325}
{
    \zA{众人都笑道：“这却使不得。”}
}
{
    \eA{for I'll go home." "This will never do!" One and all smilingly protested.}
}
{
}

\Verse{326}
{
    \zA{鸳鸯喝令小丫头子们：“拉上席去！”}
}
{
    \eA{Yüan Yang shouted to the young waiting-maids to drag her back to her table;}
}
{
}

\Verse{327}
{
    \zA{小丫头子们 也笑着，果然拉入席中。}
}
{
    \eA{and the maids, while also indulging in laughter, actually pulled her and
        compelled her to rejoin the banquet.}
}
{
}

\Verse{328}
{
    \zA{刘老老只叫：“饶了我罢！”}
}
{
    \eA{"Spare me!" old goody Liu kept on crying, "spare me!" "Any one who says one
        word more,"}
}
{
}

\Verse{329}
{
    \zA{鸳鸯道：“再多言的罚一壶。”}
}
{
    \eA{Yüan Yang exclaimed, "will be fined a whole decanter full."}
}
{
}
\ChapterImage{img/chapters/098第40回 王鳳姐擺飯秋爽齋 金鴛鴦三宣牙牌令.jpg}


\Verse{330}
{
    \zA{刘老老方住了。}
}
{
    \eA{Old goody Liu then at length observed silence.}
}
{
}

\Verse{331}
{
    \zA{鸳鸯道：“如今我说骨牌副儿，从老太太起，顺领下去，至刘老老止。}
}
{
    \eA{"I'll now give out the set of dominoes." Yüan Yang proceeded. "I'll begin
        from our venerable mistress and follow down in proper order until I come to
        old goody Liu,}
}
{
}

\Verse{332}
{
    \zA{比如我 说一副儿，将这三张牌拆开，先说头一张，再说第二张，说完了，合成这一副儿的 名字，无论诗词歌赋，成语俗话，比上一句，都要合韵。}
}
{
    \eA{when I shall stop. So as to illustrate what I meant just now by giving out a
        set, I'll take these three dominoes and place them apart; you have to begin
        by saying something on the first, next, to allude to the second, and, after
        finishing with all three,}
}
{
}

\Verse{333}
{
    \zA{错了的罚一杯。”}
}
{
    \eA{to take the name of the whole set and match it with a line,}
}
{
}

\Verse{334}
{
    \zA{众人笑 道：“这个令好，就说出来。”}
}
{
    \eA{no matter whether it be from some stanza or roundelay, song or idyl, set
        phrases or proverbs.}
}
{
}

\Verse{335}
{
    \zA{鸳鸯道：“有了一副了。}
}
{
    \eA{But they must rhyme. And any one making a mistake will be mulcted in one
        cup."}
}
{
}

\Verse{336}
{
    \zA{左边是张天。”}
}
{
    \eA{"This rule is splendid; begin at once!"}
}
{
}

\Verse{337}
{
    \zA{贾母道：“头上有青天。”}
}
{
    \eA{they all exclaimed. "I've got a set," Yüan Yang pursued; "on the left,}
}
{
}

\Verse{338}
{
    \zA{众人道好。}
}
{
    \eA{is the piece 'heaven,' (twelve dots)."}
}
{
}

\Verse{339}
{
    \zA{鸳鸯道：“当中是个五合六。”}
}
{
    \eA{"Above head stretches the blue heaven," dowager lady Chia said. "Good!"}
}
{
}

\Verse{340}
{
    \zA{贾母道：“六桥梅花香彻骨。”}
}
{
    \eA{shouted every one. "In the centre is a five and six," Yüan Yang resumed. The
        fragrance of the plum blossom pierces the bones on the bridge "Six,"}
}
{
}

\Verse{341}
{
    \zA{鸳鸯道：“剩了一 张六合么。”}
}
{
    \eA{old lady Chia added. "There now remains," Yüan Yang explained, "one piece,}
}
{
}

\Verse{342}
{
    \zA{贾母道：“一轮红日出云霄。”}
}
{
    \eA{the six and one." "From among the fleecy clouds issues the wheel-like russet
        sun."}
}
{
}

\Verse{343}
{
    \zA{鸳鸯道：“凑成却是个‘蓬头鬼’。”}
}
{
    \eA{dowager lady Chia continued. "The whole combined," Yuan Yang observed "forms
        'the devil with dishevelled hair.'"}
}
{
}

\Verse{344}
{
    \zA{贾母道：“这鬼抱住钟馗腿。”}
}
{
    \eA{"This devil clasps the leg of the 'Chung Pa' devil," old lady Chia observed.
        At the conclusion of her recitation,}
}
{
}

\Verse{345}
{
    \zA{说完，大家笑着喝彩。}
}
{
    \eA{they all burst out laughing. "Capital!" they shouted.}
}
{
}

\Verse{346}
{
    \zA{贾母饮了一杯。}
}
{
    \eA{Old lady Chia drained a cup. Yüan Yang then went on to remark,}
}
{
}

\Verse{347}
{
    \zA{鸳鸯又道：“又有一副了。}
}
{
    \eA{"I've got another set; the one on the left is a double five."}
}
{
}

\Verse{348}
{
    \zA{左边是个大长五。”}
}
{
    \eA{"Bud after bud of the plum bloom dances in the wind,"}
}
{
}

\Verse{349}
{
    \zA{薛姨妈道：“梅花朵朵风前舞。”}
}
{
    \eA{Mrs. Hsüeh replied. "The one on the right is a ten spot,"}
}
{
}

\Verse{350}
{
    \zA{鸳鸯道：“右边是个大五长。”}
}
{
    \eA{Yüan Yang pursued. "In the tenth moon the plum bloom on the hills emits its
        fragrant smell,"}
}
{
}

\Verse{351}
{
    \zA{薛姨妈道：“十月梅花岭上香。”}
}
{
    \eA{Mrs. Hsüeh added. "The middle piece is the two and five, making the 'unlike
        seven;'"}
}
{
}

\Verse{352}
{
    \zA{鸳鸯道：“当中 二五是杂七。”}
}
{
    \eA{Yüan Yang observed. "The 'spinning damsel' star meets the 'cow-herd' on the
        eve of the seventh day of the seventh moon,"}
}
{
}

\Verse{353}
{
    \zA{薛姨妈道：“织女牛郎会七夕。”}
}
{
    \eA{Miss Hsüeh said. "Together they form: 'Erh Lang strolls on the five
        mounds;'"}
}
{
}

\Verse{354}
{
    \zA{鸳鸯道：“凑成‘二郎游五岳’。”}
}
{
    \eA{Yüan Yang continued. "Mortals cannot be happy as immortals," Mrs. Hsüeh
        rejoined. Her answers over,}
}
{
}

\Verse{355}
{
    \zA{薛姨妈道：“世人不及神仙乐。”}
}
{
    \eA{the whole company extolled them and had a drink. "I've got another set!"
        Yüan Yang once more exclaimed.}
}
{
}

\Verse{356}
{
    \zA{说完，大家称赏，饮了酒。}
}
{
    \eA{"On the left, are distinctly the distant dots of the double ace."}
}
{
}

\Verse{357}
{
    \zA{鸳鸯又道：“有了一副了。}
}
{
    \eA{"Both sun and moon are so suspended as to shine on heaven and earth,"}
}
{
}

\Verse{358}
{
    \zA{左边长么两点明。”}
}
{
    \eA{Hsiang-yün ventured. "On the right,}
}
{
}

\Verse{359}
{
    \zA{湘云道：“双悬日月照乾坤。”}
}
{
    \eA{are a couple of spots, far apart, which clearly form a one and one."}
}
{
}

\Verse{360}
{
    \zA{鸳鸯道：“右边长么两点明。”}
}
{
    \eA{Yüan Yang pursued. "What time a lonesome flower falls to the ground,}
}
{
}

\Verse{361}
{
    \zA{湘云道：“闲花落地听无声。”}
}
{
    \eA{no sound is audible," Hsiang-yün rejoined. "In the middle, there is the one
        and four,"}
}
{
}

\Verse{362}
{
    \zA{鸳鸯道：“中间还 得么四来。”}
}
{
    \eA{Yüan Yang added. "The red apricot tree is planted by the sun, and leans
        against the clouds;"}
}
{
}

\Verse{363}
{
    \zA{湘云道：“日边红杏倚云栽。”}
}
{
    \eA{Hsiang-yün answered. "Together they form the 'cherry fruit ripens for the
        ninth time,'"}
}
{
}

\Verse{364}
{
    \zA{鸳鸯道：“凑成一个‘樱桃九熟’。”}
}
{
    \eA{Yüan Yang said. "In the imperial garden it is pecked by birds."}
}
{
}

\Verse{365}
{
    \zA{湘云道：“御园却被鸟衔出。”}
}
{
    \eA{Hsiang-yün replied. When she had done with her part, she drank a cup of
        wine.}
}
{
}

\Verse{366}
{
    \zA{说完，饮了一杯。}
}
{
    \eA{"I've got another set," Yüan Yang began,}
}
{
}

\Verse{367}
{
    \zA{鸳鸯道：“有了一副了。}
}
{
    \eA{"the one on the left is a double three."}
}
{
}

\Verse{368}
{
    \zA{左边是长三。”}
}
{
    \eA{"The swallows, pair by pair,}
}
{
}

\Verse{369}
{
    \zA{宝钗道：“双双燕子语梁间。”}
}
{
    \eA{chatter on the beams;" Pao-ch'ai remarked. "The right piece is a six,"}
}
{
}

\Verse{370}
{
    \zA{鸳鸯 道：“右边是三长。”}
}
{
    \eA{Yüan Yang added. "The marsh flower is stretched by the breeze e'en to the
        length of a green sash,"}
}
{
}

\Verse{371}
{
    \zA{宝钗道：“水荇牵风翠带长。”}
}
{
    \eA{Pao-ch'ai returned. "The centre piece is a three and six,}
}
{
}

\Verse{372}
{
    \zA{鸳鸯道：“当中三六九点在。”}
}
{
    \eA{making a nine spot," Yüan Yang pursued. "The three hills tower half beyond
        the azure skies;"}
}
{
}

\Verse{373}
{
    \zA{宝钗道：“三山半落青天外。”}
}
{
    \eA{Pao-ch'ai rejoined. "Lumped together they form: a 'chain-bound solitary
        boat,'"}
}
{
}

\Verse{374}
{
    \zA{鸳鸯道：“凑成‘铁练锁孤舟’。”}
}
{
    \eA{Yüan Yang resumed. "Where there are wind and waves, there I feel sad;"
        Pao-ch'ai answered.}
}
{
}

\Verse{375}
{
    \zA{宝钗道：“处 处风波处处愁。”}
}
{
    \eA{When she had finished her turn and drained her cup, Yüan Yang went on again.}
}
{
}

\Verse{376}
{
    \zA{说完饮毕。}
}
{
    \eA{"On the left," she said,}
}
{
}

\Verse{377}
{
    \zA{鸳鸯又道：“左边一个天。”}
}
{
    \eA{"there's a 'heaven.'" "A morning fine and beauteous scenery,}
}
{
}

\Verse{378}
{
    \zA{黛玉道：“良辰美景奈何天。”}
}
{
    \eA{but, alas, what a day for me!" Tai-yü replied. When this line fell on
        Pao-chai's ear,}
}
{
}

\Verse{379}
{
    \zA{宝钗听了，回头 看着他，黛玉只顾怕罚，也不理论。}
}
{
    \eA{she turned her head round and cast a glance at her, but Tai-yü was so
        nervous lest she should have to pay a forfeit that she did not so much as
        notice her.}
}
{
}

\Verse{380}
{
    \zA{鸳鸯道：“中间锦屏颜色俏。”}
}
{
    \eA{"In the middle there's the 'colour of the embroidered screen, (ten spots,
        four and six),}
}
{
}

\Verse{381}
{
    \zA{黛玉道：“纱 窗也没有红娘报。”}
}
{
    \eA{is beautiful,'" Yüan Yang proceeded. "Not e'en Hung Niang to the gauze
        window comes,}
}
{
}

\Verse{382}
{
    \zA{鸳鸯道：“剩了二六八点齐。”}
}
{
    \eA{any message to bring." Tai-yü responded. "There now remains a two and six,}
}
{
}

\Verse{383}
{
    \zA{黛玉道：“双瞻玉座引朝仪。”}
}
{
    \eA{eight in all," Yüan Yang resumed. "Twice see the jady throne when led in to
        perform the court ritual,"}
}
{
}

\Verse{384}
{
    \zA{鸳鸯道：“凑成‘篮子’好采花。”}
}
{
    \eA{Tai-yü replied. "Together they form 'a basket suitable for putting plucked
        flowers in,'"}
}
{
}

\Verse{385}
{
    \zA{黛玉道：“仙杖香挑芍药花。”}
}
{
    \eA{Yüan Yang continued. "The fairy wand smells nice as on it hangs a peony."}
}
{
}

\Verse{386}
{
    \zA{说完，饮了一 口。}
}
{
    \eA{Tai-yü retorted. At the close of her replies,}
}
{
}

\Verse{387}
{
    \zA{鸳鸯道：“左边四五成花九。”}
}
{
    \eA{she took a sip of wine. Yüan Yang then resumed. "On the left," she said,}
}
{
}

\Verse{388}
{
    \zA{迎春道：“桃花带雨浓。”}
}
{
    \eA{"there's a four and five, making a 'different-combined nine.'"}
}
{
}

\Verse{389}
{
    \zA{众人笑道：“该罚! 错了韵，而且又不像。”}
}
{
    \eA{"The peach blossoms bear heavy drops of rain;" Ying Ch'un remarked. The
        company laughed.}
}
{
}

\Verse{390}
{
    \zA{迎春笑着，饮了一口。}
}
{
    \eA{"She must be fined!" they exclaimed.}
}
{
}

\Verse{391}
{
    \zA{原是凤姐和鸳鸯都要听刘老老的笑话儿，故意都叫说错了。}
}
{
    \eA{"She has made a mistake in the rhyme. Besides, it isn't right!" Ying Ch'un
        smiled and drank a sip.}
}
{
}

\Verse{392}
{
    \zA{至王夫人，鸳鸯便 代说了一个，下便该刘老老。}
}
{
    \eA{The fact is that both lady Feng and Yüan Yang were so eager to hear the
        funny things that would be uttered by old goody Liu,}
}
{
}

\Verse{393}
{
    \zA{刘老老道：“我们庄家闲了，也常会几个人弄这个儿， 可不像这么好听就是了。}
}
{
    \eA{that they with one voice purposely ruled that every one answered wrong and
        fined them. When it came to Madame Wang's turn, Yüan Yang recited something
        for her. Next followed old goody Liu.}
}
{
}

\Verse{394}
{
    \zA{少不得我也试试。”}
}
{
    \eA{"When we country-people have got nothing to do,"}
}
{
}

\Verse{395}
{
    \zA{众人都笑道：“容易的，你只管说， 不相干。”}
}
{
    \eA{old goody Liu said, "a few of us too often come together and play this sort
        of game; but the answers we give are not so high-flown;}
}
{
}

\Verse{396}
{
    \zA{鸳鸯笑道：“左边大四是个人。”}
}
{
    \eA{yet, as I can't get out of it, I'll likewise make a try!"}
}
{
}

\Verse{397}
{
    \zA{刘老老听了，想了半日，说道：“是 个庄家人罢！”}
}
{
    \eA{"It's easy enough to say what there is," one and all laughed, "so just you
        go on and don't mind!" "On the left," Yüan Yang smiled,}
}
{
}

\Verse{398}
{
    \zA{众人哄堂笑了。}
}
{
    \eA{"there's a double four, i.e. 'man.'"}
}
{
}

\Verse{399}
{
    \zA{贾母笑道：“说的好，就是这么说。”}
}
{
    \eA{Goody Liu listened intently. After considerable reflection, "It's a
        peasant!" she cried.}
}
{
}

\Verse{400}
{
    \zA{刘老老也笑 道：“我们庄家人不过是现成的本色儿，姑娘姐姐别笑。”}
}
{
    \eA{One and all in the room blurted out laughing. "Well-said!" dowager lady Chia
        observed with a laugh,}
}
{
}

\Verse{401}
{
    \zA{鸳鸯道：“中间三四绿 配红。”}
}
{
    \eA{"that's the way." "All we country-people know," old goody Liu proceeded,}
}
{
}

\Verse{402}
{
    \zA{刘老老道：“大火烧了毛毛虫。”}
}
{
    \eA{also laughing, "is just what comes within our own rough-and-ready wits,}
}
{
}

\Verse{403}
{
    \zA{众人笑道：“这是有的，还说你的本色。”}
}
{
    \eA{so young ladies and ladies pray don't poke fun at me!" "In the centre
        there's the three and four,}
}
{
}

\Verse{404}
{
    \zA{鸳鸯笑道：“右边么四真好看。”}
}
{
    \eA{green matched with red," Yüan Yang pursued. "The large fire burnt the hairy
        caterpillar;"}
}
{
}

\Verse{405}
{
    \zA{刘老老道：“一个萝卜一头蒜。”}
}
{
    \eA{old goody Liu ventured. "This will do very well!", the party laughed, "go on
        with what is in your line."}
}
{
}

\Verse{406}
{
    \zA{众人又笑了。}
}
{
    \eA{"On the right," Yüan Yang smilingly continued,}
}
{
}

\Verse{407}
{
    \zA{鸳鸯笑道：“凑成便是‘一枝花’。”}
}
{
    \eA{"there's a one and four, and is really pretty." "A turnip and a head of
        garlic."}
}
{
}

\Verse{408}
{
    \zA{刘老老两只手比着，也要笑，却又掌住了， 说道：“花儿落了结个大倭瓜。”}
}
{
    \eA{old goody Liu answered. This reply evoked further laughter from the whole
        company. "Altogether, it's a twig of flowers," Yüan Yang added laughing.
        "The flower dropped, and a huge melon formed."}
}
{
}

\Verse{409}
{
    \zA{众人听了，由不的大笑起来。}
}
{
    \eA{old goody Liu observed, while gesticulating with both her hands by way of
        illustration.}
}
{
}

\Verse{410}
{
    \zA{只听外面乱嚷嚷的，不知何事，且听下回分解。}
}
{
    \eA{The party once more exploded in loud merriment. But, reader, if you
        entertain any curiosity to hear what else was said during the banquet,}
}
{
}

\Verse{411}
{
    \zA{(本章完)}
}
{
    \eA{listen to the explanation given in the next chapter.}
}
{
}
